\documentclass[11pt,a4paper]{article}
\usepackage[utf8]{inputenc}
\usepackage[T1]{fontenc}
\usepackage[ngerman,english]{babel}
\usepackage{geometry}
\geometry{a4paper, left=2.5cm, right=2.5cm, top=2.5cm, bottom=2.5cm}
\usepackage{mathptmx}
\usepackage{helvet}
\usepackage{amsmath}
\usepackage{amssymb}
\usepackage{amsthm}
\usepackage{titlesec}
\usepackage{booktabs}
\usepackage{tabularx}
\usepackage{xcolor}
\usepackage{authblk}
\usepackage{hyperref}
\usepackage{enumitem}
\usepackage{graphicx}
\usepackage{float}
\usepackage{setspace}
\usepackage{natbib}
\usepackage{longtable}
\usepackage{quoting}
\usepackage{multirow}
\usepackage{array}

\newtheorem{proposition}{Proposition}
\newtheorem{definition}{Definition}

\titleformat{\section}{\Large\bfseries\sffamily\color{black}}{\thesection}{1em}{}
\titleformat{\subsection}{\large\bfseries\sffamily\color{darkgray}}{\thesubsection}{1em}{}
\titleformat{\subsubsection}{\normalsize\bfseries\sffamily\color{darkgray}}{\thesubsubsection}{1em}{}

\hypersetup{
    pdftitle={Multiaxiale Psychodiagnostik},
    pdfauthor={Lukas Geiger},
    colorlinks=true,
    linkcolor=black,
    urlcolor=blue,
    citecolor=black
}

\onehalfspacing

\begin{document}

% ===================================================================
% TITELSEITE
% ===================================================================
\selectlanguage{ngerman}

\title{\textbf{\huge Ein integriertes multiaxiales Modell\\zur computergest\"utzten psychiatrischen Diagnostik}\\[0.5em]
\Large Synthese von DSM-5-TR, ICD-11 und ICF\\in einem 6-Achsen-Expertensystem\\[0.3em]
\large Ein Methoden- und Design-Paper}

\author[1]{Lukas Geiger\thanks{Korrespondenz: Lukas Geiger, Bernau, Deutschland.}}
\affil[1]{Unabh\"angiger Forscher, Bernau}

\date{Februar 2026 --- Version 2 \\ \vspace{0.5em} \small \textit{Methoden- und Design-Paper}}

\maketitle

\begin{abstract}
\noindent Die vorliegende Arbeit pr\"asentiert ein neuartiges 6-Achsen-Modell zur computergest\"utzten multiaxialen psychiatrischen Diagnostik, das die kategoriale Diagnostik nach DSM-5-TR und ICD-11 mit der funktionalen Klassifikation der ICF in einem integrierten Expertensystem vereint. Das Modell adressiert systematisch die klinischen Defizite, die durch die Abschaffung des multiaxialen Systems im DSM-5 (2013) entstanden sind, und geht zugleich \"uber das historische DSM-IV-System hinaus: Achse~I (Psychische Profile) erfasst akute, chronische, remittierte und widerlegte Diagnosen einschlie\ss{}lich Behandlungsgeschichte und Abdeckungsanalyse; Achse~II (Biographie und Entwicklung) integriert dimensionale Pers\"onlichkeitsdiagnostik mittels PID-5, pr\"agende Erfahrungen und Grundkonflikte; Achse~III (Medizinische Synopse) bildet eine symmetrische Parallelstruktur zu Achse~I mit 13~Subachsen einschlie\ss{}lich beitragender medizinischer Faktoren, Verdachtsdiagnosen, integrierter Kausalanalyse und strukturierter Medikamentenanamnese; Achse~IV (Umwelt und Funktion) kombiniert ICF, WHODAS~2.0, das deutsche GdB-System, das Cultural Formulation Interview und ein strukturiertes Bezugspersonen- und Behandlungsnetzwerk; Achse~V (Integriertes Bedingungsmodell) verbindet erstmals Fallformulierung mit diagnostischer Klassifikation in einem operationalisierten 3P/4P-Schema; Achse~VI (Belegsammlung und klinische Sicherheit) implementiert eine systematische Evidenz-Matrix mit klinischem Bewertungsfeld, CAVE-Warnhinweise f\"ur kritische klinische Risiken, einen longitudinalen Symptomverlauf und ein Kontakt- und Beobachtungsprotokoll. Die Designphilosophie des Systems basiert auf drei Prinzipien: Skelett-Prinzip (Struktur ohne Pflichtfelder), lebendes Dokument (mitwachsende Akte) und Bericht als Snapshot (eingefrorene Momentaufnahme). Das multi-professionelle Zuordnungsmodell bietet eine exemplarische, nicht verpflichtende Achsenzuordnung f\"ur Psychologen, Mediziner, Sozialarbeiter und das interdisziplin\"are Team. Die 6-Stufen-Gatekeeper-Logik bildet exakt die von Michael~B. First publizierte Goldstandard-Sequenz ab. Die Abdeckungsanalyse --- die systematische Identifikation diagnostisch unerkl\"arter Symptome --- stellt nach unserem Kenntnisstand einen neuartigen Ansatz dar und wird in Version~4 durch quantitative Abdeckungsmetriken (\%-basierte Symptom-Diagnosen-Zuordnung) und einen priorisierten 3-Stufen-Untersuchungsplan (dringend/wichtig/Verlaufskontrolle) operationalisiert. Jede Diagnose erh\"alt eine strukturierte PRO/CONTRA-Evidenzbewertung mit expliziter Konfidenzsch\"atzung. Die HiTOP-Spektren werden direkt aus den Cross-Cutting-Ergebnissen berechnet. Neue Abschnitte adressieren dimensionale Integration (HiTOP/RDoC), Ethik und Datenschutz, Validierungsstrategie sowie Interoperabilit\"atsstandards (HL7 FHIR). Die technische Implementierung als hierarchische Zustandsmaschine mit 11~St\"orungsmodulen wird vorgestellt. Der vollst\"andige Quellcode ist \"offentlich verf\"ugbar unter \url{https://github.com/lukisch/multiaxial-diagnostic-system}.

\vspace{0.5em}
\noindent \textbf{Schl\"usselbegriffe:} Multiaxiale Diagnostik, DSM-5-TR, ICD-11, ICF, Expertensystem, Differenzialdiagnostik, Abdeckungsanalyse, Fallformulierung, Cross-Cutting Symptom Measures, Hierarchische Zustandsmaschine, PID-5, WHODAS~2.0, HiTOP, RDoC, Multi-professionell, HL7 FHIR, PRO/CONTRA-Evidenzbewertung, CAVE-Warnhinweise, Symptomverlauf, Lebendes Dokument, Kontaktprotokoll

\vspace{0.5em}
\noindent \textbf{Disziplinen:} Klinische Psychologie, Psychiatrie, Medizinische Informatik, Psychometrie, Rehabilitationswissenschaft, Soziale Arbeit
\end{abstract}

\newpage
\tableofcontents
\newpage

% ===================================================================
% CO-AUTOREN
% ===================================================================
\section*{Angaben zur KI-Nutzung und Methodik}
\addcontentsline{toc}{section}{Angaben zur KI-Nutzung und Methodik}

Die vorliegende Arbeit wurde unter intensiver Mitwirkung folgender KI-Systeme erstellt. Da ihre Beitr\"age \"uber blo\ss{}e Hilfestellungen hinausgingen, werden sie hier detailliert ausgewiesen:

\begin{description}[style=nextline, leftmargin=2cm]
\item[\textbf{Claude Opus 4.6} (Anthropic)] Co-Writer: Textgenese, Strukturierung, argumentative Ausarbeitung und systematisches Review mit L\"uckenanalyse.
\item[\textbf{Gemini} (Google DeepMind) \& \textbf{Copilot} (Microsoft)] Reviewer: Kritisches Lektorat, Pr\"ufung auf Konsistenz und systematische Literaturrecherche.
\end{description}

\vspace{1em}
\noindent\textit{Hinweis:} Trotz des erheblichen maschinellen Beitrags liegt die finale Verantwortung f\"ur den wissenschaftlichen Inhalt und die Interpretation der Ergebnisse beim menschlichen Autor.

\newpage

% ===================================================================
% METHODIK
% ===================================================================
\section*{Methodik}
\addcontentsline{toc}{section}{Methodik}

Diese Arbeit folgt einer Design-Science-Methodik f\"ur die Entwicklung eines klinischen Entscheidungsunterst\"utzungs-Artefakts. Der Forschungsprozess umfasste vier Phasen:

\textbf{Phase~1: Problemidentifikation und Anforderungsanalyse.} Ein strukturiertes Literaturreview zur Abschaffung des DSM-IV-Multiaxialsystems identifizierte f\"unf spezifische klinische Defizite: Verlust der strukturierten Aufforderung zur biopsychosozialen Beurteilung, geringe GAF-Interrater-Reliabilit\"at, inkonsistente Achse-IV-Nutzung, die k\"unstliche Achse-I/II-Grenze und die niedrige Adoptionsrate der Ersatz-V/Z-Codes. Anforderungen wurden aus diesen Defiziten, aus Firsts (2024) Goldstandard-Framework f\"ur Differenzialdiagnostik und aus aktuellen dimensionalen Ans\"atzen (HiTOP, RDoC, ICD-11 dimensionales Pers\"onlichkeitsmodell) abgeleitet.

\textbf{Phase~2: Iteratives Design.} Die 6-Achsen-Architektur wurde durch iterative Zyklen von Design, Expertenfeedback und Verfeinerung entwickelt. Designentscheidungen --- wie die symmetrische Achse-I/III-Struktur, die Innovation der Abdeckungsanalyse und das CAVE-Warnsystem --- wurden gegen die identifizierten Anforderungen evaluiert. Jede Iteration wurde dokumentiert und versioniert (V1--V5).

\textbf{Phase~3: Prototypische Implementierung.} Eine Python-Referenzimplementierung (ca.\ 1.850~Zeilen) unter Verwendung hierarchischer Zustandsmaschinen (\texttt{transitions}-Bibliothek) und einer Streamlit-Oberfl\"ache wurde entwickelt, um die Machbarkeit zu demonstrieren. Der Quellcode ist \"offentlich verf\"ugbar unter \url{https://github.com/lukisch/multiaxial-diagnostic-system}.

\textbf{Phase~4: KI-gest\"utzte Entwicklung und Review.} Manuskript und Systemdesign wurden mit substanzieller KI-Unterst\"utzung entwickelt (Claude Opus 4.6 als Co-Writer; Gemini und Copilot als Reviewer). Alle KI-Beitr\"age sind im Abschnitt zur KI-Nutzung dokumentiert. Die wissenschaftliche Verantwortung f\"ur alle Inhalte liegt beim menschlichen Autor.

\textbf{Umfang und Limitationen der vorliegenden Arbeit.} Dies ist ein Methoden- und Design-Paper. Es pr\"asentiert Architektur, Begr\"undung und prototypische Implementierung des 6-Achsen-Modells. Die empirische Validierung (Interrater-Reliabilit\"at, klinischer Nutzen, konvergente Validit\"at) ist in vier Phasen geplant (Abschnitt~\ref{sec:validierung}), aber noch nicht durchgef\"uhrt. Die Abdeckungsanalyse-Schwellenwerte (85\%/60\%) sind als initiale Werte auf Basis klinischer Erw\"agungen vorgeschlagen und bed\"urfen empirischer Kalibrierung in zuk\"unftigen Studien.

\newpage

% ===================================================================
% TEIL I: GRUNDLAGEN UND PROBLEMSTELLUNG
% ===================================================================

\section{Einleitung: Die diagnostische L\"ucke nach Abschaffung des multiaxialen Systems}
\label{sec:einleitung}

Die Abschaffung des multiaxialen Diagnosesystems im DSM-5 \citep{APA2013DSM5} hinterlie\ss{} eine strukturelle L\"ucke in der psychiatrischen Diagnostik. Das urspr\"ungliche, mit dem DSM-III \citep{APA1980DSM3} eingef\"uhrte System umfasste f\"unf Achsen: Achse~I (klinische St\"orungen), Achse~II (Pers\"onlichkeitsst\"orungen und geistige Behinderung), Achse~III (medizinische Krankheitsfaktoren), Achse~IV (psychosoziale und umweltbezogene Probleme in 9~Kategorien) und Achse~V (Globale Beurteilung des Funktionsniveaus, GAF-Score 0--100). Die APA eliminierte dieses System nach einem bereits 2004 initiierten Abschaffungsantrag aus folgenden Gr\"unden: geringe Interrater-Reliabilit\"at des GAF-Scores, konzeptuelle Vermengung von Symptomen und Funktionsniveau in Achse~V, inkonsistente klinische Nutzung von Achse~IV und die k\"unstliche Grenze zwischen Achse~I und~II.

Die klinische Gemeinschaft verlor dabei mehr als sie gewann. \citet{Probst2014Multiaxial} dokumentierte, dass die Eliminierung von Achse~IV \glqq die Notwendigkeit der Kontextber\"ucksichtigung nicht eliminiert\grqq{}. \citet{Kress2014Biopsychosocial} fanden, dass Kliniker nun \glqq umso wachsamer systematische Wege zur Erfassung biopsychosozialer Informationen\grqq{} finden m\"ussen --- ohne die strukturierten Aufforderungen des alten Systems. Die als Achse-IV-Ersatz eingef\"uhrten V/Z-Codes verzeichnen eine geringe Adoptionsrate: \citet{ErlichFirst2025SocialDeterminants} konstatieren, dass \glqq ein allgemeines Bewusstsein f\"ur diese Codes und die Bedeutung ihrer Nutzung zur Kommunikation sozialer Determinanten der Gesundheit nicht eingetreten ist\grqq{}.

Am kritischsten ist der Verlust der \textit{strukturierten Aufforderung zur umfassenden Beurteilung} \"uber biologische, psychologische und soziale Dimensionen hinweg --- eine Funktion, die das flache DSM-5-Listing nicht replizieren kann.

\textbf{Bestehende multiaxiale und integrative Ans\"atze.} Das vorliegende Modell baut auf mehreren grundlegenden Ans\"atzen zur multiaxialen und integrativen psychiatrischen Diagnostik auf, geht aber substanziell \"uber diese hinaus. \citet{Williams1985MultiaxialOrigins} dokumentierte die Urspr\"unge und fr\"uhen Kritiken des DSM-III-Multiaxialsystems und identifizierte sowohl dessen klinischen Nutzen als strukturierte Aufforderung zur umfassenden Beurteilung als auch seine operationalen Schw\"achen --- insbesondere den unzuverl\"assigen GAF-Score und die inkonsistente Anwendung der psychosozialen Stressoren-Achse. Das vorliegende Modell adressiert diese spezifischen Kritiken: Es ersetzt den GAF durch validierte Instrumente (WHODAS~2.0, ICF Core Sets), eliminiert die k\"unstliche Achse-I/II-Grenze und f\"uhrt die quantitative Abdeckungsanalyse als objektive Erg\"anzung zum klinischen Urteil ein.

\citet{Bolton2008FormulationBook} argumentierte, dass Fallformulierung --- der individualisierte narrative Bericht \"uber die Schwierigkeiten eines Patienten --- die kategoriale Diagnose erg\"anzen, nicht ersetzen sollte. Boltons zentrale Einsicht war, dass Klassifikationssysteme erfassen, \textit{was} ein Patient hat, w\"ahrend die Formulierung erfasst, \textit{warum} und \textit{wie}. Das vorliegende Modell operationalisiert diese Einsicht durch Achse~V (Integriertes Bedingungsmodell), die ein strukturiertes 3P/4P-Formulierungsschema mit expliziten R\"uckbez\"ugen auf die diagnostischen Achsen implementiert --- und damit die L\"ucke zwischen Boltons narrativer Formulierung und der strukturierten Klassifikation \"uberbr\"uckt, die klinische Systeme ben\"otigen.

\citet{Wakefield1992HarmfulDysfunction} schlug die einflussreiche \glqq Harmful Dysfunction\grqq{}-Analyse vor, wonach psychische St\"orungen Zust\"ande sind, in denen ein interner Mechanismus seine nat\"urliche Funktion nicht erf\"ullt \textit{und} die Dysfunktion dem Individuum Schaden zuf\"ugt. Dieser konzeptuelle Rahmen pr\"agt den Ansatz des vorliegenden Modells zur Differenzialdiagnostik: Die 6-Stufen-Gatekeeper-Logik (nach First, 2024) unterscheidet systematisch normale Variation von Dysfunktion, indem sie Evidenz sowohl f\"ur beeintr\"achtigte Mechanismen \textit{als auch} f\"ur klinisch bedeutsame Belastung oder Funktionsbeeintr\"achtigung verlangt. Die Abdeckungsanalyse operationalisiert dies weiter, indem sie quantifiziert, welche Symptome durch etablierte Diagnosen erkl\"art werden und welche unerkl\"art bleiben --- eine direkte Operationalisierung von Wakefields Forderung, genuine Dysfunktion statt blo\ss{}er statistischer Abweichung zu identifizieren.

Die vorliegende Arbeit stellt ein 6-Achsen-Modell vor, das jede spezifische Schw\"ache des historischen Systems adressiert und zugleich F\"ahigkeiten erg\"anzt, die keines der bisherigen Systeme besessen hat.

\subsection{Institutionelle Konvergenz}

Das vorgeschlagene Modell steht nicht isoliert, sondern antizipiert die Richtung, in die sich die psychiatrischen Institutionen selbst bewegen. Das 17-k\"opfige \textit{APA Future DSM Strategic Committee} entwickelt derzeit eine \glqq Roadmap\grqq{} mit vier Unterkomitees, die Dimensionalit\"at, Biomarker, Funktion/Lebensqualit\"at und sozio\"okonomische/kulturelle/umweltbezogene Determinanten explorieren. \citet{ErlichFirst2025SocialDeterminants} fordern explizit die R\"uckkehr zu einem biaxialen System zur Erfassung sozialer Determinanten der Gesundheit.


\subsection{Designphilosophie: Skelett, lebendes Dokument, maximale Freiheit}
\label{sec:designphilosophie}

Drei Designprinzipien leiten das gesamte System:

\begin{description}[style=nextline, leftmargin=1.5cm, font=\bfseries]
\item[Skelett-Prinzip] Das System fungiert als diagnostisches \textit{Skelett}: Es bietet f\"ur jede klinisch relevante Information einen definierten Ort, ohne dass irgendein Feld als Pflichtfeld deklariert wird. Kein Diagnostiker wird gezwungen, Abschnitte auszuf\"ullen, die f\"ur den aktuellen Fall irrelevant sind. Das System ist ein Korsett zur St\"utze --- nicht zur Einschr\"ankung --- des diagnostischen Prozesses.
\item[Lebendes Dokument] Die diagnostische Akte ist kein statisches Formular, sondern ein \textit{lebendes Dokument}, das mit dem klinischen Prozess w\"achst. Neue Testergebnisse, Laborberichte, Beobachtungsprotokolle und Bewertungen k\"onnen jederzeit integriert werden. Jede Achse ist offen f\"ur iterative Erg\"anzung: Ein Laborbericht wird in die Evidenz-Matrix (Achse~VI) eingetragen mit Datum, Bezeichnung, klinischer Bewertung und Achsenverweis; eine sp\"atere Fremdanamnese erg\"anzt das Kontaktprotokoll; ein neuer psychometrischer Test aktualisiert Achse~I.
\item[Bericht als Snapshot] Berichte --- ob Entlassbriefe, Gutachten oder Verlaufsberichte --- sind \textit{eingefrorene Momentaufnahmen} des lebenden Dokuments zu einem bestimmten Zeitpunkt. Das Dokument selbst lebt weiter und kann nach einem Bericht neue Informationen aufnehmen, w\"ahrend der Bericht den Stand zum Zeitpunkt seiner Erstellung repr\"asentiert.
\end{description}

Diese Philosophie unterscheidet das System fundamental von starren Formularsystemen: Es bietet Struktur, wo Struktur n\"utzlich ist, und Freiheit, wo Freiheit klinisch geboten ist.


% ===================================================================
% 2. DAS 6-ACHSEN-MODELL
% ===================================================================
\section{Architektur des 6-Achsen-Modells}
\label{sec:modell}

Das vorgeschlagene Modell integriert die kategoriale Diagnostik (DSM-5-TR/ICD-11) in ein epistemologisches Gesamtsystem. Es dient der Erfassung der diagnostischen Lebensspanne, der Kausalit\"atspr\"ufung zwischen Somatik und Psyche sowie der Dokumentation der Behandlungs- und Remissionsgeschichte. Ein zentrales Designprinzip ist die \textbf{symmetrische Parallelstruktur} zwischen Achse~I und Achse~III: Beide Achsen verf\"ugen \"uber analoge Subachsen (Verdachtsdiagnosen, Widerlegungen, Behandlungsgeschichte, Abdeckungsanalyse, Untersuchungsplan), sodass Psychologen und Mediziner mit identischen strukturellen Werkzeugen in ihren jeweiligen Dom\"anen arbeiten k\"onnen. Tabelle~\ref{tab:achsen} gibt eine \"Ubersicht.

\begin{longtable}{p{0.8cm}p{2.2cm}p{3.8cm}p{4cm}p{2cm}}
\caption{\"Ubersicht der diagnostischen Achsen des 6-Achsen-Modells} \label{tab:achsen} \\
\toprule
\textbf{Achse} & \textbf{Bezeichnung} & \textbf{Unterteilung} & \textbf{Kerninhalte} & \textbf{Typische Profession} \\
\midrule
\endfirsthead
\toprule
\textbf{Achse} & \textbf{Bezeichnung} & \textbf{Unterteilung} & \textbf{Kerninhalte} & \textbf{Typische Profession} \\
\midrule
\endhead
I & Psychische Profile & Ia: Akute St\"orungen \newline Ib: Chronische Verl\"aufe \newline Ic: Widerlegte Verdachte \newline Id: Remittierte Diagnosen \newline Ie: Remissionsfaktoren \newline If: Behandlungsgeschichte \newline Ig: Therapietreue \newline Ih: Verdachtsdiagnosen \newline Ii: Abdeckungsanalyse \newline Ij: Untersuchungsplan & DSM-5 Cross-Cutting; PHQ-9, PCL-5, ITQ; Differenzialdiagnostik; Chronifizierung; Zeitliche Verortung & Psychologe / Psychiater \\
\addlinespace
II & Biographie \& Entwicklung & Entwicklungshistorie \newline Pr\"agende Erfahrungen \newline Grundkonflikte \newline Pers\"onlichkeit (PID-5) & IQ, Bildung, Sozialisation; Lebensereignisse; OPD-Konflikte; PID-5-BF/+M & Psychologe / Sozialp\"adagoge \\
\addlinespace
III & Medizinische Synopse & IIIa: Akute med. Diagnosen \newline IIIb: Chron. Diagnosen (erkl.) \newline IIIc: Beitragende Faktoren \newline IIId: Remittierte Erkrank. \newline IIIe: Remissionsfaktoren \newline IIIf: Behandlungsgesch. \newline IIIg: Medikamentenadh. \newline IIIh: Verdachtsdiagnosen \newline IIIi: Abdeckungsanalyse \newline IIIj: Untersuchungsplan \newline IIIk: Kausalit\"atsanalyse \newline IIIl: Genetik/Familie \newline IIIm: Medikamenten\-anamnese & Biographische Medizin\-geschichte; Kausal\-analyse Soma--Psyche; erbbiologische Belastung; Pharmako\-therapie mit Wirkungs\-bewertung & Mediziner / Facharzt \\
\addlinespace
IV & Umwelt \& Funktion & Teilhabe (ICF) \newline Psychosoziale Umst\"ande \newline Bezugspersonen-Netzwerk \newline Cultural Formulation & WHODAS~2.0, GdB; Mini-ICF-APP; Stressoren; Behandlungsnetzwerk; CFI & Sozialarbeiter / Sozialp\"adagoge \\
\addlinespace
V & Bedingungsmodell & Pr\"adisponierende Fakt. \newline Ausl\"osende Faktoren \newline Aufrechterhaltende Fakt. \newline Protektive Faktoren \newline Klinische Synthese & 3P/4P-Modell; Achsen\-integration; Behandlungs\-planung; Prognose & Interdisziplin\"ares Team \\
\addlinespace
VI & Belegsammlung \& Klinische Sicherheit & Evidenz-Matrix (m. Bewertung) \newline CAVE-Warnhinweise \newline Symptomverlauf \newline Kontaktprotokoll & Dokumentenregister m. klin. Bewertung; CAVE-Alerts; Symptomverlauf; Kontakt-/Beobachtungs\-protokoll & Alle Professionen \\
\bottomrule
\end{longtable}

\subsection{Achse I: Psychische Profile --- Temporale Diagnostik}
\label{sec:achse1}

Achse~I geht weit \"uber die einfache Diagnoseliste des DSM-5 hinaus, indem sie Diagnosen in ihrer zeitlichen Dynamik erfasst. Die Unterteilung in zehn Subachsen (Ia--Ij) erm\"oglicht die Dokumentation des gesamten diagnostischen Lebenslaufs: aktuelle akute St\"orungen (Ia), chronische Verl\"aufe mit Chronifizierungsrisiko (Ib), explizit widerlegte Verdachtsdiagnosen mit differenzialdiagnostischer Begr\"undung (Ic), remittierte Diagnosen mit belegter Inaktivit\"at (Id), Remissionsfaktoren wie Therapie, Medikation oder spontane Bew\"altigung (Ie), vollst\"andige Behandlungsgeschichte einschlie\ss{}lich Wirkungen und Nebenwirkungen (If), Therapietreue aus Selbst- und Fremdperspektive (Ig), laufende Verdachtsdiagnosen als diagnostische Hypothesen (Ih), die Abdeckungsanalyse unerklärter Restsymptomatik (Ii) und strukturierte Untersuchungspl\"ane f\"ur n\"achste diagnostische Schritte (Ij).

Die Reihenfolge von Abdeckungsanalyse (Ii) vor Untersuchungsplan (Ij) folgt der klinischen Logik: Erst m\"ussen diagnostisch unerklärte Symptome systematisch identifiziert werden, bevor gezielte Untersuchungen zu deren Kl\"arung geplant werden k\"onnen. Der Untersuchungsplan ergibt sich somit direkt aus den L\"ucken der Abdeckungsanalyse.

\textbf{PRO/CONTRA-Evidenzbewertung:} Jede Diagnose erh\"alt eine strukturierte Evidenzbewertung mit expliziter Dokumentation der Befunde, die f\"ur (PRO) und gegen (CONTRA) die Diagnose sprechen, sowie eine numerische Konfidenzsch\"atzung (0--100\%). Diese Operationalisierung folgt dem Prinzip der differenzialdiagnostischen Transparenz: Der Kliniker muss nicht nur dokumentieren, \textit{dass} eine Diagnose gestellt wird, sondern \textit{warum} --- und welche Gegenargumente er erw\"agt und verworfen hat.

\textbf{Quantitative Abdeckungsanalyse (Ii):} Die Abdeckungsanalyse wird durch eine Symptom-Diagnosen-Matrix operationalisiert, in der jedes Symptom einer oder mehreren erkl\"arenden Diagnosen zugeordnet wird. F\"ur jedes Symptom wird ein Abdeckungsprozentsatz berechnet, der den Grad der diagnostischen Erkl\"arung quantifiziert. Die Gesamtabdeckungsmetrik (z.\,B.\ $\sim$88\%) gibt dem Kliniker eine quantitative R\"uckmeldung \"uber die Vollst\"andigkeit seiner diagnostischen Arbeit. Symptome werden als \textit{vollst\"andig abgedeckt} ($\geq$85\%), \textit{partiell abgedeckt} (60--84\%) oder \textit{unzureichend abgedeckt} ($<$60\%) klassifiziert.

\textbf{Priorisierter Untersuchungsplan (Ij):} Der Untersuchungsplan verwendet ein 3-Stufen-Priorisierungsschema: \textit{Dringend} (innerhalb von 4~Wochen abzukl\"aren --- z.\,B.\ Suizidalit\"atsabkl\"arung, medizinische Notfalldiagnostik), \textit{Wichtig} (innerhalb von 3~Monaten --- z.\,B.\ neuropsychologische Testung, Differenzialdiagnostik) und \textit{Verlaufskontrolle} (kontinuierliches Monitoring --- z.\,B.\ Symptomverlauf unter Therapie). Jede geplante Untersuchung wird einem Fachgebiet zugeordnet und mit einer klinischen Begr\"undung versehen.

Diese temporale Dimensionalit\"at fehlte sowohl in Achse~I des DSM-IV als auch im flachen Listing des DSM-5.

\subsection{Achse II: Biographie und dimensionale Pers\"onlichkeit}

Achse~II modernisiert die alte Pers\"onlichkeitsachse durch Integration des dimensionalen Trait-Assessments mittels PID-5, wie es sowohl ICD-11 als auch das Alternative Modell des DSM-5 empfehlen. Das ICD-11-Modell ersetzt die kategorialen Pers\"onlichkeitsst\"orungs-Typen durch ein Schweregrad-Rating (Pers\"onlichkeitsschwierigkeit $\to$ leicht $\to$ moderat $\to$ schwer), f\"unf Trait-Qualifikatoren (Negative Affektivit\"at, Distanziertheit, Dissozialit\"at, Enthemmung, Anankastie) und einen Borderline-Muster-Spezifikator. Das DSM-5-AMPD verwendet die Level of Personality Functioning Scale (LPFS) f\"ur Kriterium~A und f\"unf Trait-Dom\"anen mit 25 Facetten f\"ur Kriterium~B.

Der \textbf{PID-5-BF+M} (36~Items, 6~Dom\"anen, 18~Facetten) verbr\"uckt beide Systeme und ist in \"uber 12~Sprachen verf\"ugbar. Metaanalysen zeigen eine Gesamt-Konvergenz von $r = 0{,}62$ zwischen den Systemen, mit Dom\"anen-Korrelationen von $r = 0{,}78$--$0{,}86$, mit Ausnahme von Anankastie ($r = 0{,}34$), die kein direktes AMPD-\"Aquivalent besitzt.

\"Uber die dimensionale Pers\"onlichkeitsdiagnostik hinaus integriert Achse~II zwei weitere biographische Komponenten: \textbf{Pr\"agende Erfahrungen} dokumentieren lebensgeschichtlich bedeutsame Ereignisse mit Lebensabschnitt und Auswirkung auf die aktuelle Situation (z.\,B.\ fr\"uhe Verlusterfahrung, Gewalterfahrung, Migration). \textbf{Grundkonflikte} erfassen wiederkehrende intrapsychische Themen (z.\,B.\ Autonomie-Abh\"angigkeits-Konflikt, Selbstwertkonflikt), wie sie in der psychodynamischen Tradition und im OPD-System beschrieben werden. Beide Komponenten sind als offene Listen implementiert, die mit dem diagnostischen Prozess wachsen, und liefern wesentliches Material f\"ur die Fallformulierung in Achse~V.

\subsection{Achse III: Medizinische Synopse --- Symmetrische Parallelstruktur zu Achse~I}
\label{sec:achse3}

Ein zentrales Designprinzip des Modells ist die \textbf{strukturelle Symmetrie} zwischen Achse~I (psychisch) und Achse~III (somatisch). Die Einsicht dahinter: Wenn ein Psychologe in Achse~I Verdachtsdiagnosen formulieren, Widerlegungen dokumentieren und Abdeckungsanalysen durchf\"uhren kann, muss ein Mediziner in Achse~III \"uber exakt dieselben M\"oglichkeiten verf\"ugen. Nur so funktioniert das multi-professionelle Modell.

Achse~III umfasst daher 13~Subachsen:

\begin{description}[style=nextline, leftmargin=1.5cm, font=\bfseries]
\item[IIIa--IIIj: Parallele Kernstruktur] Diese zehn Subachsen spiegeln die Struktur von Achse~I: akute somatische Diagnosen (IIIa), chronische somatische Diagnosen mit vollst\"andig erkl\"arendem Charakter (IIIb), beitragende medizinische Faktoren --- somatische Befunde, die Psychopathologie beg\"unstigen oder verst\"arken, ohne sie vollst\"andig zu erkl\"aren (IIIc), remittierte somatische Erkrankungen (IIId), somatische Remissionsfaktoren --- operative Eingriffe, Medikation, Lebensstil\"anderung (IIIe), medizinische Behandlungsgeschichte einschlie\ss{}lich Operationen, Pharmakotherapie und Nebenwirkungen (IIIf), Medikamentenadhärenz aus \"arztlicher Dokumentation und Patientenbericht (IIIg), medizinische Verdachtsdiagnosen als laufende diagnostische Hypothesen (IIIh), somatische Abdeckungsanalyse --- K\"orpersymptome, die durch keine aktuelle medizinische Diagnose erkl\"art werden (IIIi), und medizinischer Untersuchungsplan f\"ur weiterf\"uhrende somatische Diagnostik (IIIj).
\item[IIIk--IIIm: Achse-III-spezifische Subachsen] Drei Subachsen haben keine Entsprechung in Achse~I, da sie die spezifische Soma--Psyche-Beziehung abbilden: integrierte Kausalit\"atsanalyse --- systematische Zuordnung somatischer Befunde zur Psychopathologie unter Differenzierung von vollst\"andiger Erkl\"arung (z.\,B.\ Hypothyreose erkl\"art Depression) und beg\"unstigender Wirkung (z.\,B.\ chronischer Schmerz verst\"arkt Angstst\"orung) (IIIk); genetische und famili\"are Belastung --- Familienanamnese f\"ur somatische und psychische Erkrankungen, bekannte genetische Pr\"adispositionen und ggf.\ Ergebnisse aus Exom-Sequenzierung (IIIl); Medikamentenanamnese und Interaktionen --- strukturierte Erfassung aller aktuellen und vergangenen Medikamente mit Dosierung, Einheit, Zweck/Indikation, Einnahmeschema, Wirkungsbewertung (0--10), Nebenwirkungen und potentiellen Interaktionen (IIIm).
\end{description}

\textbf{Anmerkung zur Strukturentscheidung IIIb/IIIc:} Die explizite Differenzierung zwischen vollst\"andig erkl\"arenden chronischen Diagnosen (IIIb) und blo\ss{} beitragenden Faktoren (IIIc) integriert die Kausalanalyse direkt in die Klassifikationsstruktur. Widerlegte medizinische Verdachtsdiagnosen werden innerhalb von IIIh (Verdachtsdiagnosen) durch einen Statuswechsel dokumentiert, analog zur klinischen Praxis, in der Verdachtsdiagnosen ausgeschlossen und mit Begr\"undung verworfen werden.

\textbf{Anmerkung zu IIIm:} Die Medikamentenanamnese als eigenst\"andige Subachse (statt blo\ss{}er Auflistung in IIIf) erm\"oglicht eine strukturierte Erfassung von Wirkungsprofilen, Interaktionsrisiken und Therapietreue auf Pr\"aparatebene. Dies ist klinisch essentiell, da Polypharmazie und Medikamenteninteraktionen eine h\"aufige Ursache sowohl f\"ur somatische als auch f\"ur psychiatrische Komplikationen darstellen.

Diese Symmetrie stellt sicher, dass das System als \textit{gemeinsames Werkzeug} f\"ur Mediziner und Psychologen funktioniert. Ein Mediziner, der Achse~III ausf\"ullt, verf\"ugt \"uber dieselbe diagnostische Tiefe wie ein Psychologe in Achse~I.

\subsection{Achse IV: Umwelt und Funktion --- ICF-Integration und Cultural Formulation}
\label{sec:achse4}

Achse~IV kombiniert mehrere validierte Instrumente. Der \textbf{WHODAS~2.0} \citep{WHO2010WHODAS} (WHO Disability Assessment Schedule) erfasst sechs Dom\"anen --- Kognition, Mobilit\"at, Selbstversorgung, Umgang mit Menschen, Lebensaktivit\"aten und Partizipation --- wahlweise als 12-Item-Kurzform (erkl\"art 81\% der Varianz) oder 36-Item-Vollversion. Die psychometrischen Eigenschaften sind stark: Cronbachs $\alpha = 0{,}94$--$0{,}96$, Test-Retest-ICC $= 0{,}93$--$0{,}96$.

WHODAS~2.0 und GAF messen \textit{fundamental unterschiedliche Konstrukte}: GAF vermengt Symptome und Funktionsniveau; WHODAS~2.0 misst Behinderung unabh\"angig von der Symptomschwere. \citet{Gspandl2018WHODAS} fanden \textit{keine signifikante Korrelation} zwischen selbstbewertetem WHODAS~2.0 und GAF bei Schizophrenie-Spektrum-St\"orungen. Das System implementiert beide: WHODAS~2.0 f\"ur die Behinderungserfassung und GAF als vertraute klinische Kurzformel mit explizitem Hinweis auf dessen Limitierungen.

F\"ur den deutschen Kontext ist die \textbf{GdB-Integration} (Grad der Behinderung) essentiell. Der GdB verwendet eine 20--100-Skala ($\geq 50$ = Schwerbehinderung) und wird nach den Versorgungsmedizinischen Grunds\"atzen (VMG) bewertet.

Drei validierte \textbf{ICF Core Sets} existieren f\"ur psychische Gesundheit: Depression (31~kurze/121~umfassende Kategorien), bipolare St\"orungen (19/38) und Schizophrenie (25/97). F\"ur die praktische Implementierung ist das \textbf{Mini-ICF-APP} (Mini-ICF f\"ur Aktivit\"aten und Partizipation bei psychischen St\"orungen) von \citet{LindenBaron2005MiniICF} mit 13~Kapazit\"atsdimensionen das effizienteste Werkzeug.

\textbf{Operationalisierung der ICF-Integration in Achsen~III und~VI:} Die ICF-Codierung im 6-Achsen-Modell folgt einer dreistufigen Logik. \textit{Erstens} werden auf Achse~III (Medizinische Synopse) die K\"orperfunktionen (ICF-Kapitel~b) und K\"orperstrukturen (ICF-Kapitel~s) dokumentiert, die f\"ur die somatisch-psychische Kausalanalyse (IIIk) relevant sind --- beispielsweise \texttt{b130} (Funktionen der psychischen Energie und des Antriebs), \texttt{b152} (emotionale Funktionen) oder \texttt{b144} (Funktionen des Ged\"achtnisses). \textit{Zweitens} erfasst Achse~IV die Aktivit\"aten und Partizipation (ICF-Kapitel~d) sowie Umweltfaktoren (ICF-Kapitel~e) --- operationalisiert \"uber WHODAS~2.0 (Dom\"anen d1--d6) und Mini-ICF-APP (13~Kapazit\"atsdimensionen, z.\,B.\ F\"ahigkeit zur Regelm\"a\ss{}igkeit, Flexibilit\"at, Kontaktf\"ahigkeit). \textit{Drittens} werden in der Evidenz-Matrix (Achse~VI) die ICF-Qualifikatoren (Ausma\ss{} der Beeintr\"achtigung: 0--4) zu jeder codierten Funktionseinschr\"ankung dokumentiert, um die Quantifizierung \"uber die Zeit zu erm\"oglichen.

\textbf{Mehrwert gegen\"uber dem WHODAS~2.0 allein:} W\"ahrend der WHODAS~2.0 als globales Behinderungsinstrument sechs Dom\"anen auf einer Summenskala erfasst, erm\"oglicht das 6-Achsen-Modell eine \textit{achsenbezogene Differenzierung}: Funktionseinschr\"ankungen werden nicht nur quantifiziert, sondern kausal den Achsen~I (psychisch) und~III (somatisch) zugeordnet. Die Abdeckungsanalyse (Ii/IIIi) identifiziert zudem Funktionsdefizite, die durch keine aktuelle Diagnose erkl\"art werden --- eine Dimension, die WHODAS~2.0 nicht adressiert. Das Mini-ICF-APP erg\"anzt den WHODAS~2.0 durch st\"orungsspezifische Kapazit\"atsprofile, die f\"ur die Behandlungsplanung und Prognose in Achse~V (Bedingungsmodell) direkt verwertbar sind.

\subsubsection{Bezugspersonen und Behandlungsnetzwerk}

Achse~IV erfasst neben den psychosozialen Stressoren und Funktionsinstrumenten auch das \textbf{Bezugspersonen- und Behandlungsnetzwerk} des Patienten. F\"ur jede Kontaktperson werden Name, Rolle (Elternteil, Partner, Hausarzt, Facharzt, Therapeut, Sozialarbeiter, Betreuer etc.), Institution, Kontaktdaten und Anmerkungen dokumentiert. Diese strukturierte Netzwerkerfassung dient mehreren Zwecken: Sie macht das soziale St\"utzsystem sichtbar, das f\"ur die Fallformulierung (Achse~V, protektive Faktoren) und die Behandlungsplanung essentiell ist; sie dokumentiert das beteiligte professionelle Netzwerk f\"ur die Koordination; und sie identifiziert potentielle Informationsquellen f\"ur Fremdanamnesen (dokumentiert im Kontaktprotokoll, Achse~VI).

\subsubsection{Cultural Formulation Interview (CFI)}

Das DSM-5-TR enth\"alt ein strukturiertes \textbf{Cultural Formulation Interview} (CFI) mit 16~Kernfragen in vier Dom\"anen: kulturelle Definition des Problems, kulturelle Wahrnehmung von Ursache/Kontext/Unterst\"utzung, kulturelle Faktoren in der Bew\"altigung und kulturelle Faktoren in der Kliniker-Patient-Beziehung. F\"ur ein Modell, das umfassende biopsychosoziale Erfassung beansprucht, ist die Integration des CFI als optionales Modul in Achse~IV essentiell. Die 12~supplementären Module des CFI (f\"ur spezifische Populationen und Kontexte) k\"onnen bedarfsorientiert aktiviert werden.

\subsection{Achse V: Integriertes Bedingungsmodell --- Operationalisierte Fallformulierung}
\label{sec:achse5}

Achse~V ist g\"anzlich neuartig. Das pr\"adisponierende--ausl\"osende--aufrechterhaltende Faktorenmodell (klinisches 3P/4P-Modell) verbr\"uckt diagnostische Klassifikation mit Fallformulierung --- etwas, das keine DSM-Edition jemals enthalten hat. \citet{Owen2023Formulation} argumentiert, dass dieser Formulierungsansatz \glqq die Auswahl der Fakten diszipliniert und die Behandlung zielgerichtet macht\grqq{}. Die Synthese aller Achsen --- wie Biologie (III), Biographie (II) und aktuelle Stressoren (IV) mit der Psychopathologie (I) interagieren --- wird hier expliziert.

Das Modell operationalisiert die Fallformulierung durch vier strukturierte Komponenten:

\begin{description}[style=nextline, leftmargin=1.5cm, font=\bfseries]
\item[Pr\"adisponierende Faktoren (P1)] Vulnerabilit\"atsfaktoren, die vor der St\"orung existierten. Quellen: Achse~II (biographische Risikofaktoren, Pers\"onlichkeits-Traits), Achse~III (genetische Belastung, IIIm), Achse~IV (fr\"uhe psychosoziale Belastungen). Kodierung: Faktor, Quellenachse, Evidenzst\"arke (gesichert/wahrscheinlich/m\"oglich), Beleg (Achse~VI-Referenz).
\item[Ausl\"osende Faktoren (P2)] Ereignisse oder Ver\"anderungen, die das Auftreten der St\"orung ausgelöst haben. Quellen: Achse~IV (aktuelle Stressoren, Life Events), Achse~III (somatische Ausl\"oser). Kodierung: Faktor, zeitlicher Zusammenhang, Quellenachse, Beleg.
\item[Aufrechterhaltende Faktoren (P3)] Bedingungen, die die St\"orung perpetuieren. Quellen: Achse~I (Komorbidit\"aten, Therapietreue), Achse~III (unbehandelte somatische Befunde), Achse~IV (fortbestehende psychosoziale Belastungen). Kodierung: Faktor, Mechanismus, \"Anderbarkeit (modifizierbar/stabil), Priorit\"at f\"ur Intervention.
\item[Protektive Faktoren (P4)] Ressourcen, die Resilienz f\"ordern. Quellen: Achse~II (Pers\"onlichkeitsst\"arken), Achse~IV (soziale Unterst\"utzung, \"okonomische Ressourcen), Achse~I (Remissionsfaktoren, Ie). Kodierung: Faktor, Quellenachse, Aktivierbarkeit.
\end{description}

Die Verkn\"upfung mit der Abdeckungsanalyse (Achse~Ii) ist zentral: Symptome, die durch keine Achse-I-Diagnose abgedeckt sind, \textit{sollten} durch die Achse-V-Formulierung erkl\"arbar sein. Symptome, die weder diagnostisch noch formulatorisch erkl\"art werden, repr\"asentieren genuine diagnostische L\"ucken und triggern den Untersuchungsplan (Achse~Ij).

\subsection{Achse VI: Belegsammlung, CAVE-Warnhinweise und Symptomverlauf}

Achse~VI wurde gegen\"uber dem Originalentwurf substanziell erweitert und umfasst nun drei Komponenten:

\begin{description}[style=nextline, leftmargin=1.5cm, font=\bfseries]
\item[Evidenz-Matrix] Ein zentrales Dokumentenverzeichnis, das jede kodierte Information mit einem Beleg verkn\"upft (Dokumentenanalyse, Befragung, Testung). Jeder Eintrag umfasst neben Achsenbezug, Dokumententyp und Beschreibung ein explizites \textbf{Bewertungsfeld} f\"ur die klinische Einsch\"atzung: Ein eingehender Laborbericht wird mit Datum, Bezeichnung, Quelle \textit{und} klinischer Bewertung dokumentiert --- so ist transparent, wie der Kliniker den Befund im diagnostischen Kontext interpretiert. Diese Innovation hat keinen Pr\"azedenzfall in irgendeinem Klassifikationssystem und unterst\"utzt direkt die klinische Rechenschaftspflicht.
\item[CAVE-Warnhinweise] Ein strukturiertes Warnsystem f\"ur kritische klinische Risiken, die achsen\"ubergreifend Relevanz haben. Jeder Warnhinweis wird mit einer Kategorie versehen: \textit{Medikamenten-Interaktion} (z.\,B.\ Lithium + NSAR), \textit{Labor-Artefakt} (z.\,B.\ CRP-Erh\"ohung durch chronische Entz\"undung, nicht akute Infektion), \textit{Kontraindikation} (z.\,B.\ Schwangerschaft bei bestimmter Medikation), \textit{zeitliche Fehlzuordnung} (z.\,B.\ Symptombeginn vor oder nach vermutetem Ausl\"oser), \textit{diagnostische Einschr\"ankung} (z.\,B.\ nicht verwertbarer Testbefund) und \textit{sonstiger Warnhinweis}. Jeder Alert referenziert die betroffene Quellenachse (I--VI), sodass die klinische Relevanz sofort kontextualisiert ist. CAVE-Alerts werden in der Gesamtsynopse prominent rot hervorgehoben.
\item[Longitudinaler Symptomverlauf] Eine strukturierte Dokumentation des zeitlichen Symptomverlaufs: F\"ur jedes relevante Symptom wird Beginn (Erstmanifestation), aktueller Status (aktiv/remittiert/fluktuierend) und Therapieansprechen (Response/Non-Response/Partial Response) erfasst. Diese longitudinale Perspektive erm\"oglicht die Identifikation von Verlaufsmustern, Therapieresistenz und Chronifizierungsrisiken, die in einer rein querschnittlichen Diagnostik unsichtbar bleiben.
\item[Kontakt- und Beobachtungsprotokoll] Ein strukturiertes Protokoll f\"ur alle klinisch relevanten Kontakte und Beobachtungen. Jeder Eintrag umfasst Datum, Kontaktart (Telefonat, Gespr\"ach, Beobachtung, Hausbesuch, Fremdanamnese, E-Mail/Brief), beteiligte Kontaktperson, Inhalt und optionalen Achsenbezug. Das Protokoll dokumentiert den \textit{Prozess} der Informationsgewinnung: Wann wurde mit wem gesprochen, was wurde beobachtet, und welche Achse profitiert von dieser Information? Damit fungiert Achse~VI nicht nur als statische Belegsammlung, sondern als \textbf{Prozessdokumentation} des diagnostischen Arbeitens --- ein wesentlicher Bestandteil des Konzepts des lebenden Dokuments (vgl.\ Abschnitt~\ref{sec:designphilosophie}).
\end{description}

Die Erweiterung von Achse~VI reflektiert die klinische Erkenntnis, dass ein diagnostisches System nicht nur \textit{klassifizieren}, sondern auch \textit{warnen}, \textit{\"uber die Zeit verfolgen} und den \textit{diagnostischen Prozess selbst dokumentieren} muss. Die CAVE-Funktion adressiert das in der Polypharmazie- und Komplexit\"atsmedizin wachsende Problem klinischer Risiken, die in keiner einzelnen Achse vollst\"andig abgebildet sind, aber achsen\"ubergreifend Konsequenzen haben.

\subsection{Multi-professionelle Zuordnung}
\label{sec:multiprofessionell}

Das 6-Achsen-Modell bietet eine \textbf{empfohlene}, nicht verpflichtende Zuordnung von Achsen zu Professionsgruppen. Tabelle~\ref{tab:professionen} zeigt eine exemplarische Konfiguration f\"ur interdisziplin\"are Teams. In der klinischen Praxis k\"onnen die Zust\"andigkeiten flexibel angepasst werden: Ein Psychiater, der sowohl psychodiagnostische als auch somatische Kompetenz besitzt, kann selbstverst\"andlich Achse~I \textit{und} Achse~III ausf\"ullen; eine Fach\"arztin f\"ur Psychosomatische Medizin kann alle sechs Achsen bedienen. Das System \textit{erm\"oglicht} interdisziplin\"are Arbeitsteilung, \textit{erzwingt} sie aber nicht.

\begin{table}[H]
\centering
\caption{Exemplarische multi-professionelle Achsenzuordnung (nicht verpflichtend)}
\label{tab:professionen}
\begin{tabularx}{\textwidth}{lXX}
\toprule
\textbf{Achse} & \textbf{Typische Profession} & \textbf{Fachliche Begr\"undung} \\
\midrule
I & Psychologe / Psychiater & Psychodiagnostische Kompetenz; Testungshoheit \\
\addlinespace
II & Psychologe / Sozialp\"adagoge & Biographische Anamnese; Pers\"onlichkeitsdiagnostik \\
\addlinespace
III & Mediziner / Facharzt & Somatische Diagnosestellung; Kausalanalyse \\
\addlinespace
IV & Sozialarbeiter / Sozialp\"adagoge & Lebensweltexpertise; ICF-Kompetenz; Teilhabebewertung \\
\addlinespace
V & Interdisziplin\"ares Team & Synthese profitiert von allen Perspektiven \\
\addlinespace
VI & Alle Professionen & Jede beteiligte Profession belegt ihre eigenen Befunde \\
\bottomrule
\end{tabularx}
\end{table}

Der Mehrwert der Zuordnung liegt nicht in einer starren Zugangsbeschr\"ankung, sondern in der \textbf{strukturierten Aufforderung}: Das System macht transparent, \textit{welche} Kompetenz f\"ur welche Achse typischerweise ben\"otigt wird. In Einzelpraxen oder kleineren Settings kann eine einzelne Fachperson alle Achsen bearbeiten; in Kliniken und interdisziplin\"aren Teams erm\"oglicht das Modell eine effiziente Arbeitsteilung, bei der jede Profession in ihrer Dom\"ane mit identischen strukturellen Werkzeugen arbeitet. Die identische Subachsen-Struktur von Achse~I und~III ist dabei kein Zufall, sondern Designentscheidung: Ein Mediziner, der in IIIh eine somatische Verdachtsdiagnose formuliert, nutzt exakt dieselbe Logik wie ein Psychologe, der in Ih eine psychische Verdachtsdiagnose eintr\"agt.

\subsection{Komorbidit\"atsregeln und diagnostische Hierarchie}
\label{sec:komorbiditaet}

Das System implementiert drei Ebenen von Komorbidit\"atsregeln, die f\"ur die korrekte diagnostische Entscheidungsfindung essentiell sind:

\begin{description}[style=nextline, leftmargin=1.5cm, font=\bfseries]
\item[Hierarchische Exklusionsregeln] Bestimmte Diagnosen schlie\ss{}en andere aus. Beispiel: Eine Schizophrenie-Diagnose schlie\ss{}t eine gleichzeitige schizoaffektive St\"orung aus. Das System kodiert diese als harte Constraints, die bei der Diagnoseeingabe automatisch gepr\"uft werden.
\item[\glqq Due to another condition\grqq{}-Regeln] Viele DSM-5-TR-Diagnosen erfordern den Ausschluss, dass die Symptomatik durch eine andere psychische St\"orung, eine Substanz oder einen medizinischen Krankheitsfaktor besser erkl\"art wird. Diese Regeln werden direkt durch die Gatekeeper-Stufen~1--3 (Abschnitt~\ref{sec:gatekeeper}) implementiert.
\item[Erlaubte Komorbidit\"aten mit Priorisierung] Bei 4+ aktiven Diagnosen priorisiert das System nach: (1)~klinischer Dringlichkeit (Suizidalit\"at, Psychose), (2)~Schweregrad (Achse-I-Akutstatus), (3)~Behandelbarkeit (modifizierbare aufrechterhaltende Faktoren aus Achse~V), (4)~chronologischer Reihenfolge.
\end{description}


% ===================================================================
% 3. GATEKEEPER-LOGIK
% ===================================================================
\section{Die 6-Stufen-Gatekeeper-Logik der Differenzialdiagnostik}
\label{sec:gatekeeper}

Das System implementiert eine dynamische Warteschlange, in der Screening-Auff\"alligkeiten spezifische DSM-5-TR/ICD-11-Module triggern. Die 6-Stufen-Sequenz bildet exakt das von Michael~B. First im \textit{DSM-5-TR Handbook of Differential Diagnosis} (2024) publizierte Framework ab (Tabelle~\ref{tab:gatekeeper}).

\begin{table}[H]
\centering
\caption{Alignment der 6-Stufen-Gatekeeper-Logik mit Firsts Goldstandard}
\label{tab:gatekeeper}
\begin{tabularx}{\textwidth}{clXc}
\toprule
\textbf{Stufe} & \textbf{Systemschritt} & \textbf{Firsts Framework} & \textbf{Match} \\
\midrule
1 & Simulationsausschluss & Malingering und artifizielle St\"orung ausschlie\ss{}en & Exakt \\
\addlinespace
2 & Substanzausschluss & Substanz\"atiologie ausschlie\ss{}en & Exakt \\
\addlinespace
3 & Medizinischer Ausschluss & \"Atiologische medizinische Erkrankung ausschlie\ss{}en & Exakt \\
\addlinespace
4 & Prim\"arkategorie & Spezifische Prim\"arst\"orung(en) bestimmen & Exakt \\
\addlinespace
5 & Anpassungsst\"orung & Anpassungsst\"orungen differenzieren & Exakt \\
\addlinespace
6 & Funktionsschwelle & Grenze zu \glqq keine psychische St\"orung\grqq{} & Konzeptuell \\
\bottomrule
\end{tabularx}
\end{table}

\citet{First2024Handbook} beschreibt dieses Framework als den definitiven Makro-Level-Diagnoseprozess. Sein Handbook enth\"alt dar\"uber hinaus \textbf{30 symptomorientierte Entscheidungsb\"aume} (zwei in der TR-Edition erg\"anzt f\"ur dissoziative Symptome und repetitive pathologische Verhaltensweisen) und \textbf{67~Differenzialdiagnosetabellen}, die alle dieser Sequenz folgen.

\textbf{Anmerkung zu Stufe~6 (\glqq konzeptuell\grqq{}):} W\"ahrend die Stufen~1--5 exakt algorithmisch abbildbar sind, ist Stufe~6 --- die Grenzziehung zwischen normaler Stressreaktion und psychischer St\"orung --- inh\"arent eine klinische Urteilsentscheidung. Das System unterst\"utzt diese durch objektive Funktionsdaten (WHODAS~2.0, Mini-ICF-APP aus Achse~IV) und die Cross-Cutting-Schwellenwerte, kann aber die klinische Entscheidung nicht vollst\"andig automatisieren. Diese epistemische Grenze wird transparent kommuniziert.


% ===================================================================
% 4. CROSS-CUTTING SCREENING
% ===================================================================
\section{Cross-Cutting Symptom Measures als intelligente Triage}
\label{sec:screening}

Die DSM-5 Cross-Cutting Symptom Measures bilden das R\"uckgrat der Screening-Architektur. Das \textbf{Level-1-Erwachsenenma\ss{}} enth\"alt \textbf{23~Items \"uber 13~Dom\"anen}: Depression (2), \"Arger (1), Manie (2), Angst (3), somatische Symptome (2), Suizidalit\"at (1), Psychose (2), Schlafprobleme (1), Ged\"achtnis (1), repetitive Gedanken und Verhaltensweisen (2), Dissoziation (1), Pers\"onlichkeitsfunktion (2) und Substanzgebrauch (3). Jedes Item verwendet eine 5-stufige Likert-Skala (0~=~keine bis 4~=~schwer).

Die Schwellenlogik ist klinisch kalibriert: \textbf{Die meisten Dom\"anen triggern Level~2 bei $\geq 2$ (leicht)}, aber drei sicherheitskritische Dom\"anen --- Suizidalit\"at, Psychose und Substanzgebrauch --- \textbf{triggern bei $\geq 1$ (gering)}. Diese Asymmetrie reflektiert das klinische Imperativ, selbst minimale Best\"atigung gef\"ahrlicher Symptome nie zu \"ubersehen.

Die Level-2-Ma\ss{}e verweisen auf spezifische validierte Instrumente: PROMIS Depression Short Form, PROMIS Anxiety Short Form, Altman Self-Rating Mania Scale (ASRM), PHQ-15 f\"ur somatische Symptome, PROMIS Sleep Disturbance, adaptierte FOCI Severity Scale f\"ur Zwangsst\"orungen und adaptierter NIDA-Modified ASSIST f\"ur Substanzgebrauch. F\"unf Dom\"anen (Suizidalit\"at, Psychose, Ged\"achtnis, Dissoziation, Pers\"onlichkeitsfunktion) besitzen \textit{keine offiziellen Level-2-Ma\ss{}e} --- die APA empfiehlt hier klinische Evaluation.

Die DSM-5 Field Trials \citep{Narrow2013FieldTrials} zeigten \textbf{gute bis exzellente Test-Retest-Reliabilit\"at} (ICC~0,64--0,97) f\"ur die meisten Items.

\subsection{Kinder- und Jugendversion}

Das DSM-5 Cross-Cutting-System umfasst eine separate \textbf{Elternberichtsversion f\"ur Kinder und Jugendliche} (6--17~Jahre) mit \textbf{25~Items \"uber 12~Dom\"anen}. Die Dom\"ane \glqq Pers\"onlichkeitsfunktion\grqq{} entf\"allt entwicklungsbedingt; stattdessen werden Reizbarkeit und \"Arger st\"arker gewichtet. F\"ur eine vollst\"andige Implementierung sollte das System die p\"adiatrische Variante als separaten Einstiegspfad anbieten und altersangepasste Level-2-Instrumente referenzieren (z.\,B.\ CBCL, SDQ, SCARED).

\subsection{Umfassende Screening-Instrument-Matrix}

Tabelle~\ref{tab:instrumente} zeigt die empfohlenen Instrumente f\"ur alle wesentlichen diagnostischen Dom\"anen.

\begin{longtable}{p{2.5cm}p{2.3cm}cp{1.5cm}p{2.5cm}c}
\caption{Screening-Instrument-Matrix} \label{tab:instrumente} \\
\toprule
\textbf{Dom\"ane} & \textbf{Instrument} & \textbf{Items} & \textbf{Cutoff} & \textbf{Sens./Spez.} & \textbf{Kosten} \\
\midrule
\endfirsthead
\toprule
\textbf{Dom\"ane} & \textbf{Instrument} & \textbf{Items} & \textbf{Cutoff} & \textbf{Sens./Spez.} & \textbf{Kosten} \\
\midrule
\endhead
Depression & PHQ-9 & 9 & $\geq 10$ & 88\%/88\% & Frei \\
\addlinespace
Angst & GAD-7 & 7 & $\geq 10$ & 89\%/82\% & Frei \\
\addlinespace
PTBS (DSM-5) & PCL-5 & 20 & 31--33 & 85--95\%/82--90\% & Frei \\
\addlinespace
PTBS/KPTBS (ICD-11) & ITQ & 18 & Algorithmus & Exzellent & Frei \\
\addlinespace
Psychoserisiko & PQ-16 & 16 & $\geq 6$ & 87\%/87\% & Frei \\
\addlinespace
Bipolar-Screening & MDQ & 15 & $\geq 7$ & 73\%/90\% & Frei \\
\addlinespace
ADHS & ASRS~v1.1 & 6 & $\geq 4$ & 69\%/99,5\% & Frei \\
\addlinespace
Autismus & AQ-10 & 10 & $\geq 6$ & 88\%/91\% & Frei \\
\addlinespace
Alkoholgebrauch & AUDIT & 10 & $\geq 8$ & 92\%/94\% & Frei \\
\addlinespace
Drogengebrauch & DAST-10 & 10 & $\geq 3$ & 98\%/91\% & Frei \\
\addlinespace
Pers\"onlichkeit & PID-5-BF & 25 & Dimensional & N/A & Frei \\
\addlinespace
Suizidalit\"at & C-SSRS & 6 & Jede Best\"at. & Risikoklassif. & Frei \\
\addlinespace
Zwang & OCI-R & 18 & $\geq 21$ & Gut & Frei \\
\addlinespace
Somatisierung & SSS-8 & 8 & $\geq 12$ & Vergleichbar & Frei \\
\addlinespace
Dissoziation & DES-II & 28 & $\geq 30$ & 74\%/80\% & Frei \\
\addlinespace
Essst\"orungen & SCOFF & 5 & $\geq 2$ & 84\%/90\% & Frei \\
\addlinespace
Essst\"orungen (detail.) & EDE-QS & 12 & $\geq 15$ & Gut/Gut & Frei \\
\addlinespace
Schlafst\"orungen & ISI & 7 & $\geq 15$ & 82\%/82\% & Frei \\
\addlinespace
Impulskontrolle & SSIS & 5 & $\geq 3$ & Screening & Frei \\
\bottomrule
\end{longtable}

Alle empfohlenen Instrumente sind \textbf{frei verf\"ugbar} --- keine Lizenzkosten behindern die Implementierung.


% ===================================================================
% 5. KLASSIFIKATIONSDIVERGENZEN
% ===================================================================
\section{Kritische Divergenzen zwischen DSM-5-TR und ICD-11}
\label{sec:divergenzen}

Das System muss f\"unf kritische Divergenzen zwischen den Klassifikationssystemen kodieren.

\subsection{Schizophrenie-Dauer}

DSM-5-TR (295.90/F20.x) verlangt \textbf{6~Monate} kontinuierlicher Zeichen einschlie\ss{}lich mindestens 1~Monat aktiver Symptome, w\"ahrend ICD-11 (6A20) nur \textbf{1~Monat} fordert. Ein Patient kann somit die ICD-11-Kriterien f\"ur Schizophrenie erf\"ullen, unter DSM-5 aber nur die Diagnose einer schizophreniformen St\"orung erhalten.

\subsection{PTBS-Architektur}

DSM-5-TR verwendet \textbf{4~Cluster mit 20~Symptomen} (Intrusion $\geq$1/5, Vermeidung $\geq$1/2, Negative Kognitionen $\geq$2/7, Arousal $\geq$2/6). ICD-11 verwendet ein bewusst schmaleres Modell: \textbf{3~Cluster mit 6~Kernsymptomen}. ICD-11 f\"ugt dann die \textbf{Komplexe PTBS (6B41)} als distinkte Diagnose hinzu, die alle PTBS-Kriterien plus drei St\"orungen der Selbstorganisation (Affektdysregulation, negatives Selbstkonzept, Beziehungsschwierigkeiten) erfordert.

\subsection{Pers\"onlichkeitsst\"orungen}

ICD-11 ersetzte kategoriale Typen vollst\"andig durch ein dimensionales Modell. Das DSM-5-AMPD verwendet die LPFS f\"ur Kriterium~A und f\"unf Trait-Dom\"anen mit 25~Facetten f\"ur Kriterium~B. Der PID-5-BF+M verbr\"uckt beide Systeme.

\subsection{Gaming Disorder}

ICD-11 (6C51) verwendet einen monothetischen Ansatz (alle 4~Kriterien erforderlich; $\geq 12$~Monate). DSM-5-TR listet Internet Gaming Disorder nur als \glqq Condition for Further Study\grqq{} mit polythetischem Ansatz ($\geq 5$ von 9~Kriterien). Konkordanz: $\kappa = 0{,}80$, aber ICD-11 hat eine h\"ohere diagnostische Schwelle (Pr\"avalenz 2,7\% vs.\ DSM-5 5,2\%).

\subsection{Komplexe Trauer / Prolonged Grief Disorder}

ICD-11 (6B42) und DSM-5-TR (Prolonged Grief Disorder, neu aufgenommen) kodieren Anh\"altende Trauer\"st\"orung, unterscheiden sich aber in Zeitkriterium (ICD-11: 6~Monate; DSM-5-TR: 12~Monate) und Symptomkonfiguration. Das System muss beide Varianten abbilden.


% ===================================================================
% 6. ABDECKUNGSANALYSE
% ===================================================================
\section{Die Abdeckungsanalyse: Ein vorgeschlagener Ansatz}
\label{sec:abdeckung}

Die Abdeckungsanalyse (\textit{Coverage Analysis}) stellt die innovativste Komponente des Systems dar. Die Forschungsliteratur best\"atigt, dass \textbf{kein publiziertes Werkzeug eine automatisierte Abdeckungsanalyse implementiert} --- die systematische Kennzeichnung von Symptomen, die durch aktuelle Diagnosen nicht erkl\"art werden.

Die engsten existierenden Parallelen sind: Komorbidit\"ats-Detektionsalgorithmen, die Symptome kennzeichnen, die auf zus\"atzliche Diagnosen hindeuten; Firsts Differenzialdiagnosetabellen \citep{First2024Handbook}, die \"uberlappende Pr\"asentationen vergleichen; und transdiagnostische Frameworks wie HiTOP, in denen Symptome als Netzwerkknoten modelliert werden. \citet{Nordgaard2020Transdiagnostic} zeigten empirisch, dass die meisten Symptome \"uber die meisten St\"orungen hinweg auftreten --- \textbf{Depressions- und Angstsymptome fanden sich bei fast allen Erstaufnahme-Patienten unabh\"angig von der Diagnose} --- was den Bedarf an einem systematischen Abdeckungsanalyse-Werkzeug direkt st\"utzt.

Die Implementierung arbeitet wie folgt: (1)~Sammlung aller best\"atigten Symptome \"uber alle Screening-Instrumente, (2)~Zuordnung jedes Symptoms zu den diagnostischen Kriterien, die es erf\"ullt, (3)~Berechnung eines Abdeckungsprozentsatzes pro Symptom, der den Grad der diagnostischen Erkl\"arung quantifiziert, (4)~Klassifikation als vollst\"andig abgedeckt ($\geq$85\%), partiell abgedeckt (60--84\%) oder unzureichend abgedeckt ($<$60\%), (5)~Berechnung einer Gesamtabdeckungsmetrik als gewichteter Mittelwert \"uber alle Symptome, (6)~Pr\"asentation der Ergebnisse als Symptom-Diagnosen-Matrix mit visueller Hervorhebung diagnostischer L\"ucken. Das 4P-Modell (Achse~V) dient als nat\"urliches Komplement: Symptome, die nicht durch Achse-I-Diagnosen abgedeckt sind (Subachse~Ii), sollten durch die Achse-V-Formulierung erkl\"arbar sein; solche, die es nicht sind, repr\"asentieren genuine diagnostische L\"ucken und m\"unden in den priorisierten Untersuchungsplan (Subachse~Ij).

Die quantitative Metrik (z.\,B.\ \glqq Gesamtabdeckung: $\sim$88\%\grqq{}) liefert dem Kliniker eine unmittelbare R\"uckmeldung \"uber die Vollst\"andigkeit seiner diagnostischen Arbeit --- ein Feature, das in der gesamten publizierten psychiatrischen Informatik ohne Pr\"azedenzfall ist.

Die Abdeckungsanalyse wird symmetrisch auch in Achse~III implementiert (Subachse~IIIi): K\"orpersymptome, die durch keine aktuelle somatische Diagnose erkl\"art werden, werden systematisch identifiziert und f\"uhren zum medizinischen Untersuchungsplan (Subachse~IIIj).


% ===================================================================
% 7. DIMENSIONALE INTEGRATION
% ===================================================================
\section{Dimensionale Integration: HiTOP und RDoC als erg\"anzende Perspektiven}
\label{sec:dimensional}

Das 6-Achsen-Modell basiert prim\"ar auf kategorialer Diagnostik (DSM-5-TR/ICD-11), enth\"alt aber bereits dimensionale Elemente (PID-5, WHODAS~2.0). Zwei weitere dimensionale Frameworks verdienen Integration als erg\"anzende Schichten:

\subsection{HiTOP (Hierarchical Taxonomy of Psychopathology)}

Die Hierarchical Taxonomy of Psychopathology \citep{Kotov2017HiTOP} organisiert psychische St\"orungen empirisch-hierarchisch in sechs Spektren: Internalizing (Depression, Angst, PTBS), Thought Disorder (Psychose, Manie), Disinhibited Externalizing (Substanzgebrauch, Impulskontrolle), Antagonistic Externalizing (Antisoziales Verhalten), Detachment (soziale Zur\"uckgezogenheit, Anhedonie) und Somatoform (somatische Symptome). Die Cross-Cutting-Ergebnisse des Systems werden direkt auf HiTOP-Spektren gemappt und als Radardiagramm visualisiert. Das Mapping nutzt die maximale Auspr\"agung der zugeh\"origen Cross-Cutting-Dom\"anen:

\begin{itemize}[noitemsep]
\item Internalizing = max(Depression, Anxiety, Somatic, Sleep)
\item Thought Disorder = max(Psychosis, Dissociation)
\item Disinhibited Externalizing = max(Substance, Mania)
\item Antagonistic Externalizing = max(Anger)
\item Detachment = max(Detachment/Memory)
\item Somatoform = max(Somatic)
\end{itemize}

Diese automatische Berechnung erfordert \textit{keine zus\"atzliche Datenerhebung} --- die bereits im Cross-Cutting-Screening (Abschnitt~\ref{sec:screening}) erfassten Dom\"anenwerte werden direkt wiederverwendet. Die Darstellung als Radardiagramm (visuell distinct vom PID-5-Pers\"onlichkeitsprofil) f\"ordert die Erkennung transdiagnostischer Muster, die bei rein kategorialer Diagnostik unsichtbar bleiben.

\subsection{RDoC (Research Domain Criteria)}

Die Research Domain Criteria des NIMH \citep{Insel2010RDoC} definieren sechs Dom\"anen (Negative Valence, Positive Valence, Cognitive Systems, Social Processes, Arousal/Regulatory, Sensorimotor) mit sieben Analyseebenen (Gene, Molek\"ule, Zellen, Schaltkreise, Physiologie, Verhalten, Selbstbericht). F\"ur ein klinisches System ist RDoC prim\"ar als \textit{Forschungsannotation} relevant: Wenn Biomarker-Daten verf\"ugbar sind (z.\,B.\ Cortisol-Profile, Neuroimaging), k\"onnen diese den RDoC-Dom\"anen zugeordnet werden. Das APA Future DSM Strategic Committee exploriert explizit die Integration von Biomarkern --- das System ist darauf vorbereitet.

Ein separates Konzeptpapier beschreibt die detaillierte Architektur der dimensionalen Integration.


% ===================================================================
% 8. TECHNISCHE ARCHITEKTUR
% ===================================================================
\section{Technische Architektur der Python-Implementierung}
\label{sec:technik}

\subsection{Hierarchische Zustandsmaschinen als Entscheidungsmotor}

Nach Evaluation von vier Architekturans\"atzen --- AnyTree, Zustandsmaschinen, Rule Engines und Behavior Trees --- erweisen sich \textbf{Hierarchische Zustandsmaschinen (HSMs)} mittels der Python-Bibliothek \texttt{transitions} als optimale L\"osung f\"ur psychiatrische Diagnostik-Workflows.

Die \texttt{transitions}-Bibliothek mit ihrer \texttt{HierarchicalMachine}-Erweiterung bietet: konditionelle Transitionen via Guards (direkte Abbildung von \glqq wenn Symptom~X $\to$ betrete Modul~Y\grqq{}), verschachtelte Zust\"ande (der 6-Stufen-Prozess mit st\"orungsspezifischen Submodulen), History States (R\"ucknavigation) und serialisierbaren Zustand (Speichern/Fortsetzen).

Die Architektur bildet sich wie folgt ab:

\begin{verbatim}
Top-level: [Intake -> Step1_Malingering -> Step2_Substance ->
            Step3_Medical -> Step4_CrossCutting ->
            Step5_DisorderModules -> Step6_Functioning -> Summary]

Step5_DisorderModules (verschachtelt, 11 Module):
  +-- MoodDisorders -> {MDD_Criteria, Bipolar_Screening, Dysthymia,
                        PMDD, Prolonged_Grief}
  +-- AnxietyDisorders -> {GAD, Panic, Social_Anxiety, Phobias,
                           Separation_Anxiety, Selective_Mutism}
  +-- TraumaDisorders -> {PTSD_DSM5, PTSD_ICD11, CPTSD_DSO,
                          Acute_Stress, Adjustment}
  +-- PsychoticDisorders -> {Schizophrenia_1mo, Schizophrenia_6mo,
                             Schizoaffective, Brief_Psychotic}
  +-- PersonalityDisorders -> {LPFS, PID5_Traits, ICD11_Severity}
  +-- OCDSpectrum -> {OCD_Criteria, BDD, Hoarding, Trichotillomania,
                      Excoriation}
  +-- DissociativeDisorders -> {DID, Depersonalization,
                                Dissociative_Amnesia}
  +-- EatingDisorders -> {AN, BN, BED, ARFID}
  +-- NeurodevelopmentalDisorders -> {ADHD, ASD_Assessment}
  +-- SubstanceUseDisorders -> {Alcohol_Use, Drug_Use,
                                Behavioral_Addictions}
  +-- SomaticDisorders -> {Somatic_Symptom, Illness_Anxiety,
                           Conversion, Factitious}
\end{verbatim}

Die Erweiterung von 5 auf \textbf{11~St\"orungsmodule} stellt sicher, dass jedes Screening-Instrument der Matrix (Tabelle~\ref{tab:instrumente}) in einen diagnostischen Workflow m\"undet. Ein auff\"alliger DES-II-Score triggert das Modul DissociativeDisorders; ein erh\"ohter SCOFF-Score das Modul EatingDisorders. Ohne diese Vollst\"andigkeit bestünde eine systemische L\"ucke zwischen Screening und Diagnostik.

AnyTree wird weiterhin f\"ur \textbf{Visualisierung und Serialisierung} der Baumstruktur genutzt (Graphviz-Export, JSON-Repr\"asentation), w\"ahrend \texttt{transitions} die Laufzeitlogik verwaltet.

\subsection{Empfohlener Technologie-Stack}

\begin{table}[H]
\centering
\caption{Empfohlener Technologie-Stack}
\label{tab:techstack}
\begin{tabularx}{\textwidth}{lXl}
\toprule
\textbf{Komponente} & \textbf{Empfehlung} & \textbf{Begr\"undung} \\
\midrule
Entscheidungsmotor & \texttt{transitions} (HSM) & Verschachtelte Zust\"ande, Guards \\
Baumvisualisierung & \texttt{anytree} & Graphviz-Export, JSON \\
UI (Prototyp) & Streamlit & Schnelles Prototyping, \texttt{st.navigation} \\
PDF-Berichte & WeasyPrint + Jinja2 & Natives UTF-8, HTML/CSS-Templates \\
Pers\"onlichkeits-Charts & Plotly & \texttt{px.line\_polar()} f\"ur PID-5-Radardiagramme \\
Datenpersistenz & SQLModel + SQLite & Pydantic + SQLAlchemy kombiniert \\
Datenvalidierung & Pydantic & Typsichere Diagnosekriterien-Modelle \\
ICD-11-Codes & WHO ICD-11 API & Strukturierte Diagnose-Code-Suche \\
Interoperabilit\"at & HL7 FHIR (R4) & Standardisierter Datenaustausch \\
\bottomrule
\end{tabularx}
\end{table}

\subsection{Open-Source-Referenzimplementierung}

Der vollst\"andige Quellcode des Prototyps (V9, ca.\ 1.850~Zeilen Python/Streamlit), die bilinguale \"Ubersetzungsdatei, der Entwicklungsfahrplan und das vorliegende Paper sind \"offentlich verf\"ugbar unter:

\begin{center}
\url{https://github.com/lukisch/multiaxial-diagnostic-system}
\end{center}

\noindent Die Bereitstellung als Open-Source-Repository dient der wissenschaftlichen Transparenz und Reproduzierbarkeit. Die Implementierung kann mit \texttt{pip install -r requirements.txt} und \texttt{streamlit run multiaxial\_diagnostic\_system.py} unmittelbar gestartet werden.

\textbf{Lizenzierung und Archivierung:} Das Repository steht unter der MIT-Lizenz, die sowohl akademische als auch kommerzielle Nutzung erlaubt, jedoch keinerlei Gew\"ahrleistung \"ubernimmt (vgl.\ Haftungsfragen, Abschnitt~\ref{sec:ethik}). F\"ur die langfristige Zitierbarkeit und Versionierung sollte das Repository mittels Zenodo archiviert werden. Eine DOI-Registrierung (z.\,B.\ \texttt{10.5281/zenodo.XXXXXX}) erm\"oglicht die dauerhafte Referenzierung einer spezifischen Version, unabh\"angig von \"Anderungen am GitHub-Repository. \textit{Hinweis:} Zum Zeitpunkt der Einreichung dieses Papers ist die Zenodo-Archivierung noch ausstehend und sollte vor Publikation nachgeholt werden.


% ===================================================================
% 9. FUNKTIONALE ASSESSMENT-INTEGRATION
% ===================================================================
\section{Funktionale Beurteilung: Br\"ucke zwischen Diagnose und Behinderung}
\label{sec:funktional}

Die Integration funktionaler Beurteilung verbindet Diagnose und Teilhabe. Das Modell implementiert drei Ebenen: das krankheitsspezifische Assessment (ICF Core Sets), die allgemeine Behinderungsmessung (WHODAS~2.0) und die rechtlich-administrative Bewertung (GdB).

ICD-11 selbst hat sich in Richtung ICF-Integration bewegt, indem \glqq Functioning Properties\grqq{} eingef\"uhrt wurden, die mit 103 rehabilitationsrelevanten Gesundheitszust\"anden verkn\"upft sind. F\"ur Schizophrenie und psychotische St\"orungen enth\"alt ICD-11 \textbf{dimensionale Qualifikatoren}, die mit rehabilitationsbasierter psychiatrischer Versorgung konsistent sind.


% ===================================================================
% 10. ETHIK UND DATENSCHUTZ
% ===================================================================
\section{Ethik, Datenschutz und algorithmischer Bias}
\label{sec:ethik}

Die computergest\"utzte psychiatrische Diagnostik wirft spezifische ethische Fragen auf, die bei der Entwicklung und Implementierung systematisch adressiert werden m\"ussen.

\subsection{Algorithmischer Bias und Validierungspopulationen}

Screening-Instrumente sind an spezifischen Populationen validiert, und die \"Ubertragbarkeit ihrer Cutoff-Werte auf andere Populationen ist limitiert. Der PHQ-9 zeigt kulturell variable Cutoff-Werte: W\"ahrend f\"ur westliche Populationen ein Cutoff von $\geq 10$ validiert ist, liegen in asiatischen L\"andern h\"aufig niedrigere optimale Schwellenwerte vor. Der AQ-10 hat eine bekannte Geschlechterverzerrung bei Frauen mit Autismus-Spektrum-St\"orungen (\"Uberdiagnose bei M\"annern, Unterdiagnose bei Frauen aufgrund unterschiedlicher Ph\"anotypen). Das AUDIT (Alkoholscreening) untersch\"atzt systematisch Alkoholprobleme bei Frauen, da die Cutoffs an m\"annlichen Trinkmustern kalibriert wurden.

\textbf{Sozio\"okonomischer Bias:} Instrumente wie das WHODAS~2.0 erfassen Funktionseinschr\"ankungen, die bei Personen mit niedrigem sozio\"okonomischem Status teilweise durch fehlende Ressourcen (z.\,B.\ Transport, Kinderbetreuung) bedingt sind, nicht durch die psychische St\"orung selbst. Das System muss diese strukturellen Faktoren in Achse~IV (psychosoziale Umst\"ande) separat dokumentieren, um Fehlattributionen zu vermeiden.

\textbf{Migrationshintergrund und Sprachbarrieren:} Selbstberichts-Instrumente setzen ausreichende Sprachkompetenz voraus. Bei Personen mit Migrationshintergrund k\"onnen sprachliche Missverst\"andnisse zu falsch-positiven oder falsch-negativen Ergebnissen f\"uhren. Das System sollte bei nicht-muttersprachlichen Patienten einen expliziten Hinweis auf diese Limitation ausgeben.

\textbf{Mitigation:} Das System muss diese Bias-Risiken transparent dokumentieren und --- wo verf\"ugbar --- populationsspezifische Normen anbieten. Die Integration des Cultural Formulation Interview (Abschnitt~\ref{sec:achse4}) adressiert kulturellen Bias auf der Systemebene. F\"ur jedes Screening-Instrument sollte die Validierungspopulation dokumentiert sein, sodass Kliniker die \"Ubertragbarkeit auf den konkreten Fall beurteilen k\"onnen.

\subsection{Datenschutz und DSGVO}

Psychiatrische Diagnosedaten geh\"oren zu den sensibelsten Gesundheitsdaten. Die Implementierung muss folgende Anforderungen erf\"ullen: Verschl\"usselung aller gespeicherten und \"ubertragenen Daten (AES-256, TLS~1.3); optionale Zugriffskontrolle nach Professionszuordnung (Abschnitt~\ref{sec:multiprofessionell}); Audit-Trail f\"ur alle diagnostischen Entscheidungen; Recht auf L\"oschung und Datenportabilit\"at; Datensparsamkeit --- nur diagnostisch notwendige Informationen werden erfasst. Besondere Vorsicht gilt f\"ur die Belegsammlung (Achse~VI), die pers\"onliche Dokumente referenziert.

\subsection{Informed Consent und Patientenautonomie}

Die Nutzung eines computergest\"utzten Diagnosesystems wirft Fragen der informierten Einwilligung auf. Patienten haben das Recht zu wissen, dass ihre diagnostische Beurteilung teilweise durch algorithmische Verfahren unterst\"utzt wird. Das System sollte daher folgende Transparenzanforderungen erf\"ullen:

\begin{itemize}[noitemsep]
\item \textbf{Aufkl\"arung \"uber Systemnutzung:} Patienten werden dar\"uber informiert, dass Screening-Instrumente und diagnostische Entscheidungsb\"aume als Unterst\"utzungswerkzeuge eingesetzt werden.
\item \textbf{Opt-out-M\"oglichkeit:} Patienten haben das Recht, eine rein menschlich durchgef\"uhrte Diagnostik zu verlangen, falls sie der computergest\"utzten Methode nicht vertrauen.
\item \textbf{Erkl\"arbarkeit von Diagnosevorschl\"agen:} Das System muss nachvollziehbar machen, welche Symptome und Kriterien zu einem Diagnosevorschlag gef\"uhrt haben. Die PRO/CONTRA-Evidenzbewertung (Abschnitt~\ref{sec:achse1}) dient genau diesem Zweck.
\item \textbf{Zugang zu eigenen Daten:} Patienten haben das Recht auf Einsicht in ihre diagnostische Akte (Achse~I--VI) und auf Erkl\"arung der dokumentierten Befunde.
\end{itemize}

Diese Anforderungen sind nicht nur ethisch geboten, sondern auch rechtlich durch die DSGVO (Art.~13--15) verankert.

\subsection{Haftungsfragen und medizinisch-rechtliche Verantwortung}

Die Nutzung eines computergest\"utzten Diagnosesystems verschiebt nicht die Haftung vom Kliniker auf das System. Folgende rechtliche Prinzipien gelten:

\begin{description}[style=nextline, leftmargin=1.5cm, font=\bfseries]
\item[Behandler-Haftung bleibt bestehen] Der behandelnde Arzt oder Psychotherapeut bleibt vollumf\"anglich haftbar f\"ur diagnostische Fehlentscheidungen, auch wenn diese durch ein algorithmisches System unterst\"utzt wurden. Das System ist ein Werkzeug, kein Entscheidungstr\"ager.
\item[Sorgfaltspflicht bei Systemnutzung] Kliniker m\"ussen das System sachgem\"a\ss{} nutzen und dessen Grenzen kennen. Eine blinde \"Ubernahme algorithmischer Diagnosevorschl\"age ohne klinische Pr\"ufung erf\"ullt nicht die \"arztliche Sorgfaltspflicht.
\item[Dokumentationspflicht] Die Nutzung des Systems und die klinische \"Uberpr\"ufung der Systemvorschl\"age sollten dokumentiert werden (Achse~VI, Kontaktprotokoll). Dies dient der Nachvollziehbarkeit im Haftungsfall.
\item[Produkthaftung des Entwicklers] Der Systementwickler haftet f\"ur nachweisliche Softwarefehler (z.\,B.\ falsche Implementierung von DSM-5-TR-Kriterien, Rechenfehler in der Abdeckungsanalyse). Die Open-Source-Ver\"offentlichung (Abschnitt~\ref{sec:technik}) erm\"oglicht externe Code-Reviews und erh\"oht die Transparenz.
\item[Keine Zulassung als Medizinprodukt erforderlich?] Nach EU-Medizinprodukteverordnung (MDR 2017/745) sind reine Dokumentationssysteme ohne therapeutische oder pr\"adiktive Funktion nicht zwingend als Medizinprodukt einzustufen. Das vorliegende System unterst\"utzt die diagnostische Dokumentation, trifft aber keine autonomen therapeutischen Entscheidungen. Eine juristische Pr\"ufung dieser Einstufung ist dennoch ratsam.
\end{description}

Die klare Trennung zwischen \textit{diagnostischer Unterst\"utzung} (zul\"assig) und \textit{autonomer Diagnosestellung} (nicht implementiert) ist essentiell f\"ur die rechtliche Absicherung.

\subsection{Klinische Verantwortung und epistemische Grenzen}

Das System ist als \textit{Unterst\"utzungswerkzeug} konzipiert, nicht als diagnostischer Automat. Alle algorithmischen Diagnosevorschl\"age erfordern klinische Best\"atigung. Die epistemische Grenze der Automatisierbarkeit (vgl.\ Stufe~6 der Gatekeeper-Logik, Abschnitt~\ref{sec:gatekeeper}) wird im System explizit kommuniziert. Die finale diagnostische Verantwortung liegt stets beim behandelnden Kliniker.

\textbf{Kritische F\"alle, die algorithmische Systeme nicht zuverl\"assig bewerten k\"onnen:}
\begin{itemize}[noitemsep]
\item \textbf{Suizidalit\"at:} Selbstberichts-Instrumente (C-SSRS) sind hilfreich, aber jede Best\"atigung suizidaler Gedanken erfordert ein ausf\"uhrliches klinisches Gespr\"ach. Algorithmische Risikobewertungen k\"onnen die klinische Einsch\"atzung nicht ersetzen.
\item \textbf{Psychotische Symptome:} Selbstberichts-Instrumente (PQ-16) sind Screening-Werkzeuge, keine diagnostischen Instrumente. Eine Best\"atigung psychotischer Symptome erfordert klinische Exploration und oft Fremdanamnese.
\item \textbf{Kulturgebundene Syndrome:} Das DSM-5-TR enth\"alt Glossare kulturgebundener Syndrome (z.\,B.\ Ataque de nervios, Taijin kyofusho), die algorithmisch schwer zu erfassen sind. Hier ist kulturelle Kompetenz des Klinikers unersetzlich.
\end{itemize}

Diese Grenzen machen das System nicht unbrauchbar, sondern unterstreichen die Notwendigkeit einer \textit{hybriden} Diagnostik, in der algorithmische Unterst\"utzung und klinisches Urteil komplementär wirken.


% ===================================================================
% 11. VALIDIERUNGSSTRATEGIE
% ===================================================================
\section{Validierungsstrategie}
\label{sec:validierung}

\textbf{Aktueller Status:} Das vorliegende Paper pr\"asentiert die konzeptionelle Architektur und technische Implementierung des 6-Achsen-Modells. Die nachfolgend beschriebene Validierungsstrategie stellt einen \textit{Implementierungsplan} dar --- die empirische Validierung ist zum gegenw\"artigen Zeitpunkt noch nicht durchgef\"uhrt. Dieser Abschnitt beschreibt daher die geplanten Validierungsschritte, nicht bereits vorliegende Ergebnisse.

Die Abschaffung des DSM-IV-Systems wurde unter anderem durch mangelnde Interrater-Reliabilit\"at begr\"undet. Das vorliegende System muss daher von Anfang an eine rigorose Validierungsstrategie verfolgen:

\begin{description}[style=nextline, leftmargin=1.5cm, font=\bfseries]
\item[Phase~1: Experten-Review (N=10--15)] Strukturiertes Review der Achsenarchitektur durch Psychiater, klinische Psychologen und Sozialarbeiter. Bewertung der Vollst\"andigkeit, klinischen Plausibilit\"at und Praktikabilit\"at. Ziel: Identifikation fehlender Subachsen oder \"uberfl\"ussiger Komponenten.

\item[Phase~2: Inter-Rater-Reliabilitätsstudie mit Fallvignetten (N=10 Rater, 20 Vignetten)] \textit{Dies ist die kritische Pilotstudie zur initialen Validierung des Systems:}

\textbf{Rekrutierung:} $N = 10$ unabh\"angige Rater (5 Psychiater, 5 approbierte Psychologische Psychotherapeuten; alle mit $\geq 3$ Jahren klinischer Erfahrung).

\textbf{Stimulusmaterial:} 20 standardisierte Fallvignetten (10 einfache F\"alle mit Einzeldiagnose, 10 komplexe F\"alle mit Komorbidit\"at), erstellt nach Vorbild der DSM-5-TR Clinical Cases \citep{BarnhillDSM5ClinicalCases}. Jede Vignette enth\"alt: Anamnese, aktuelle Beschwerden, Symptomverlauf, relevante medizinische Vorgeschichte, soziales Umfeld. Diagnosen vorab durch Expertenpanel (2 unabh\"angige Fach\"arzte + Konsensprozess) als Goldstandard etabliert.

\textbf{Prozedur:} Jeder Rater diagnostiziert alle 20 Vignetten (1) mit konventionellem Vorgehen (Freitext-Diagnose mit Begr\"undung), (2) mit dem 6-Achsen-System (vollst\"andige Achsendiagnostik). Reihenfolge der Methoden randomisiert; mindestens 1 Woche Washout zwischen Bedingungen pro Vignette.

\textbf{Prim\"are Outcome-Variablen:}
\begin{itemize}[leftmargin=*, noitemsep]
    \item \textit{Inter-Rater-Reliabilit\"at:} Fleiss' $\kappa$ f\"ur kategoriale Diagnosen (Achse~I Primärdiagnose); Intraklassen-Korrelation (ICC) f\"ur dimensionale Ratings (PID-5-BF auf Achse~II, WHODAS~2.0 auf Achse~IV). Akzeptanzkriterium: $\kappa \geq 0{,}7$ (substanzielle \"Ubereinstimmung); ICC $\geq 0{,}75$.
    \item \textit{Konvergente Validit\"at:} \"Ubereinstimmung 6-Achsen-Diagnosen vs.\ Goldstandard. Sensitivit\"at und Spezifit\"at pro Diagnose.
    \item \textit{Inkrementelle Validit\"at der Abdeckungsanalyse:} Anteil der F\"alle, in denen Achse~Ii (unerkl\"arte Symptome) klinisch relevante Differenzialdiagnosen identifiziert, die im Freitext-Vorgehen \"ubersehen wurden. Hypothese: $\geq 20$\,\% der komplexen F\"alle profitieren von Abdeckungsanalyse.
    \item \textit{Time-on-Task:} Durchschnittliche Bearbeitungszeit pro Vignette (konventionell vs.\ 6-Achsen).
\end{itemize}

\textbf{Sekund\"are Outcome-Variablen:}
\begin{itemize}[leftmargin=*, noitemsep]
    \item System Usability Scale (SUS) nach Abschluss aller Vignetten
    \item Subjektive Einsch\"atzung der Nützlichkeit einzelner Achsen (5-Punkt-Likert)
    \item Qualitative Interviews zu identifizierten Schwachstellen des Systems
\end{itemize}

\textbf{Power-Analyse:} Bei $N = 10$ Ratern und 20 Vignetten ($200$ diagnostische Urteile) kann ein ICC $= 0{,}75$ mit Power $1 - \beta = 0{,}8$ und $\alpha = 0{,}05$ detektiert werden (95\%-KI-Breite $\pm 0{,}12$).

\textbf{Falsifikationskriterium:} Falls $\kappa < 0{,}6$ oder ICC $< 0{,}65$, ist die Inter-Rater-Reliabilit\"at unzureichend --- dies erfordert Revision der Achsenstruktur oder klarere Operationalisierungen.

\item[Phase~3: Klinische Felderprobung (N=200+)] Prospektive Anwendung in klinischen Settings (ambulant und station\"ar). Messung von: Interrater-Reliabilit\"at, konvergente Validit\"at gegen\"uber etablierten Systemen, inkrementelle Validit\"at der Abdeckungsanalyse (werden durch Ii/IIIi identifizierte L\"ucken klinisch best\"atigt?), Akzeptanz und Usability (System Usability Scale, SUS).
\item[Phase~4: Multi-Site-Trial] Multizentrische Studie zur Generalisierbarkeit. Vergleich station\"ar vs.\ ambulant, Erwachsene vs.\ Kinder/Jugendliche, verschiedene kulturelle Kontexte.
\end{description}


% ===================================================================
% 12. INTEROPERABILITÄT
% ===================================================================
\section{Interoperabilit\"at und Gesundheits-IT-Standards}
\label{sec:interop}

F\"ur die Integration in bestehende Gesundheitsinformationssysteme muss das 6-Achsen-Modell standardkonforme Schnittstellen implementieren.

\textbf{HL7 FHIR (R4):} Die Achsen lassen sich auf FHIR-Ressourcen abbilden: Achse~I/III-Diagnosen als \texttt{Condition}-Ressourcen mit Extensions f\"ur Subachsen-Zuordnung; Screening-Ergebnisse als \texttt{QuestionnaireResponse}; Funktionsstatus (Achse~IV) als \texttt{ClinicalImpression}; die Fallformulierung (Achse~V) als \texttt{CarePlan}; Belege (Achse~VI) als \texttt{DocumentReference}.

\textbf{SNOMED CT:} Alle Diagnosen sollten neben DSM-5-TR- und ICD-11-Codes auch SNOMED-CT-Konzept-IDs f\"uhren, um internationale Interoperabilit\"at zu gew\"ahrleisten.

\textbf{MHIRA-Kompatibilit\"at:} Das MHIRA-Projekt (Mental Health Information Reporting Assistant, BMC Psychiatry 2023) stellt die relevanteste Open-Source-Referenzarchitektur dar: ein cloud-basiertes, Docker-deployiertes psychiatrisches EHR-System mit digitalisierten psychometrischen Instrumenten. Eine FHIR-basierte Schnittstelle zum MHIRA-System sollte priorisiert werden.


% ===================================================================
% 13. DISKUSSION UND AUSBLICK
% ===================================================================
\section{Diskussion und Ausblick}
\label{sec:diskussion}

Das vorgestellte 6-Achsen-Modell ist kein R\"uckgriff auf das DSM-IV, sondern repr\"asentiert eine \textbf{vorgeschlagene Weiterentwicklung}, auf die die psychiatrische Fachwelt selbst konvergiert. Neun Merkmale sind besonders innovativ:

\textbf{Erstens} hat die Abdeckungsanalyse (Abschnitt~\ref{sec:abdeckung}) keine existierende Implementierung in der publizierten psychiatrischen Informatik. Sie adressiert das von Nordgaard et al.\ (2016) dokumentierte Grundproblem transdiagnostischer Symptom\"uberlappung und wird symmetrisch in Achse~I (Ii) und Achse~III (IIIi) implementiert.

\textbf{Zweitens} verbr\"uckt die operationalisierte Integration von Fallformulierung (Achse~V) mit diagnostischer Klassifikation die langj\"ahrige Kluft zwischen \glqq Welche St\"orung hat dieser Patient?\grqq{} und \glqq Warum hat dieser Patient diese St\"orung?\grqq{}.

\textbf{Drittens} adressiert die Dual-System-Architektur DSM-5-TR/ICD-11 ein reales klinisches Bed\"urfnis in Kontexten, in denen beide Systeme Anwendung finden.

\textbf{Viertens} stellt die symmetrische Parallelstruktur von Achse~I und Achse~III sicher, dass das System als gemeinsames Werkzeug f\"ur alle beteiligten Professionen funktioniert --- ein Designprinzip, das in keinem bisherigen System verfolgt wurde.

\textbf{F\"unftens} bietet die Erweiterung auf 11~St\"orungsmodule erstmals eine vollst\"andige Abdeckung zwischen Screening und Diagnostik: Jedes Screening-Instrument m\"undet in einen strukturierten diagnostischen Workflow.

\textbf{Sechstens} operationalisiert die PRO/CONTRA-Evidenzbewertung mit numerischer Konfidenz die differenzialdiagnostische Transparenz: Der Kliniker muss nicht nur \textit{welche} Diagnose er stellt, sondern \textit{warum} dokumentieren --- einschlie\ss{}lich der bewusst verworfenen Gegenargumente.

\textbf{Siebtens} implementiert das CAVE-Warnsystem in Achse~VI eine achsen\"ubergreifende Sicherheitsschicht, die kritische klinische Risiken (Labor-Artefakte, Medikamenten-Interaktionen, Kontraindikationen, zeitliche Fehlzuordnungen) prominent kennzeichnet --- ein Feature, das in keinem publizierten Klassifikationssystem existiert.

\textbf{Achtens} erg\"anzt der longitudinale Symptomverlauf die querschnittliche Diagnostik um eine zeitliche Dimension: Die Dokumentation von Symptombeginn, aktuellem Status und Therapieansprechen erm\"oglicht die Identifikation von Chronifizierungsmustern und Therapieresistenz, die f\"ur die Behandlungsplanung essentiell sind.

\textbf{Neuntens} realisiert die Designphilosophie des \textit{lebenden Dokuments} ein Paradigma, das die diagnostische Akte nicht als einmalig auszuf\"ullendes Formular, sondern als mitwachsende Prozessdokumentation versteht. Bezugspersonen-Netzwerk (Achse~IV), pr\"agende Erfahrungen und Grundkonflikte (Achse~II) sowie das Kontakt- und Beobachtungsprotokoll (Achse~VI) sind offene Listen, die den diagnostischen Prozess selbst dokumentieren und kontinuierlich Material f\"ur die Fallformulierung (Achse~V) liefern.

Die technische Implementierung sollte vier Meilensteine priorisieren: (1)~das Cross-Cutting-Screening-Modul als Einstiegspunkt des Systems, (2)~die Gatekeeper-Zustandsmaschine mittels \texttt{transitions} HSM, (3)~das PID-5-BF-Assessment mit Plotly-Radardiagramm und WeasyPrint-PDF-Export, und (4)~die FHIR-Schnittstellendefinition f\"ur Interoperabilit\"at.


% ===================================================================
% 14. LIMITATIONEN UND ZUKÜNFTIGER FORSCHUNGSBEDARF
% ===================================================================
\section{Limitationen und zuk\"unftiger Forschungsbedarf}
\label{sec:limitationen}

Das vorgestellte 6-Achsen-Modell unterliegt mehreren Limitationen, die bei der Interpretation und Weiterentwicklung ber\"ucksichtigt werden m\"ussen.

\textbf{Fehlende empirische Validierung:} Die zentrale Limitation besteht darin, dass das Modell bislang rein konzeptionell-theoretisch entwickelt wurde. Weder Interrater-Reliabilit\"at noch konvergente oder inkrementelle Validit\"at sind empirisch gepr\"uft. Die in Abschnitt~\ref{sec:validierung} beschriebene 4-Phasen-Validierungsstrategie stellt einen Forschungsplan dar, keine vorliegenden Ergebnisse. Bis zur Durchf\"uhrung mindestens der Phasen~1 und~2 bleibt der klinische Nutzen des Modells eine begr\"undete Hypothese.

\textbf{Komplexit\"at und klinische Praktikabilit\"at:} Das System umfasst sechs Achsen mit insgesamt 23+ Subachsen. Obwohl das Skelett-Prinzip keine Pflichtfelder vorsieht, k\"onnte die Komplexit\"at in der klinischen Praxis abschreckend wirken. Eine wesentliche Forschungsfrage ist, ob Kliniker das System tats\"achlich in der intendierten Tiefe nutzen oder auf wenige Achsen reduzieren. Usability-Studien (SUS-Skala) sind essenziell.

\textbf{Keine Normierung der Abdeckungsanalyse:} Die quantitative Abdeckungsanalyse (Achse~Ii/IIIi) verwendet prozentuale Metriken, f\"ur die keine empirischen Normen existieren. Die Schwellenwerte (85\%/60\%) sind klinisch plausibel, aber nicht empirisch validiert. Zuk\"unftige Forschung muss kl\"aren, welche Abdeckungsquoten in verschiedenen diagnostischen Populationen typisch sind.

\textbf{Formale Definition der Abdeckungsanalyse erforderlich:} Die Abdeckungsanalyse --- als \glqq ohne publizierten Pr\"azedenzfall\grqq{} beschrieben --- verf\"ugt \"uber keine formale mathematische Definition. Die prozentbasierte Symptom-Diagnosen-Zuordnung ist intuitiv ansprechend, wirft jedoch mehrere ungel\"oste Fragen auf: Wie werden Symptomgewichte zugewiesen? Wie werden \"uberlappende Diagnosen in der Abdeckungsberechnung behandelt? Wie sensitiv ist die Gesamtabdeckungsmetrik gegen\"uber einzelnen Symptombewertungen? Eine formale Spezifikation mit durchgerechneten Beispielen ist erforderlich, bevor die Abdeckungsanalyse unabh\"angig implementiert oder evaluiert werden kann.

\textbf{Eingeschr\"ankte Literaturbasis:} Als Designarbeit st\"utzt sich diese Arbeit prim\"ar auf diagnostische Klassifikationshandb\"ucher und eine kleine Anzahl grundlegender Publikationen. Ein systematisches Review aller existierenden computergest\"utzten Diagnosesysteme (z.\,B.\ OPCRIT, DTREE, CATEGO) und deren empirische Erfolgsbilanz w\"urde die Begr\"undung f\"ur die hier getroffenen spezifischen Designentscheidungen st\"arken.

\textbf{KI-intensive Entwicklung:} Die substanzielle KI-Beteiligung sowohl am Systemdesign als auch an der Manuskripterstellung (Disclosure Level~5) wirft Fragen zur Reproduzierbarkeit und unabh\"angigen Verifizierbarkeit der Designentscheidungen auf. W\"ahrend der Quellcode \"offentlich verf\"ugbar ist, ist der iterative Designprozess --- einschlie\ss{}lich der Expertenfeedback-Zyklen --- nicht vollst\"andig so dokumentiert, dass eine unabh\"angige Replikation m\"oglich w\"are.

\textbf{Kulturelle Generalisierbarkeit:} Das Modell integriert das Cultural Formulation Interview (CFI), wurde aber prim\"ar aus westlich-europ\"aischer Perspektive entwickelt. Die Anwendbarkeit in nicht-westlichen klinischen Kontexten bedarf gesonderter Pr\"ufung. Insbesondere die Pers\"onlichkeitsdiagnostik mittels PID-5 und die biographischen Komponenten (Achse~II) k\"onnen kulturell unterschiedlich interpretiert werden.

\textbf{Technologieabh\"angigkeit:} Die Implementierung als computergest\"utztes Expertensystem setzt eine funktionsf\"ahige IT-Infrastruktur voraus. In ressourcenarmen Settings oder bei Systemausf\"allen muss das Modell auch als papierbasiertes Schema funktionieren --- eine Anforderung, die bislang nicht systematisch gepr\"uft wurde.

\textbf{Zuk\"unftiger Forschungsbedarf:} Neben der empirischen Validierung (Abschnitt~\ref{sec:validierung}) sind drei Forschungsrichtungen priorit\"ar: (1)~die Entwicklung eines abgek\"urzten \glqq Rapid Assessment\grqq{}-Modus f\"ur Notaufnahmen und Erstgespräche, (2)~die Integration pr\"adiktiver Modelle (maschinelles Lernen) zur Unterst\"utzung der Differenzialdiagnostik auf Basis der akkumulierten Patientendaten, und (3)~l\"angsschnittliche Studien zur Sensitivit\"at des Symptomverlaufs (Achse~VI) f\"ur Therapieansprechen und R\"uckfallpr\"adiktion.


% ===================================================================
% 14a. KLINISCHE FALLVIGNETTE
% ===================================================================
\subsection{Illustrative Fallvignette}
\label{sec:fallvignette}

\textit{Hinweis: Die folgende Fallvignette ist fiktiv und dient ausschlie\ss{}lich der Illustration der Anwendbarkeit des 6-Achsen-Modells. Jede \"Ahnlichkeit mit realen Personen ist zuf\"allig.}

\vspace{0.5em}
\noindent \textbf{Patient M., 34~Jahre, m\"annlich, Erstvorstellung in der psychiatrischen Ambulanz.}

\begin{description}[style=nextline, leftmargin=1.5cm, font=\bfseries\small]
\item[Achse~I (Psychische Profile)]
Ia (akut): Mittelgradige depressive Episode (F32.1), Konfidenz 85\%, PRO: PHQ-9 = 16, Antriebslosigkeit, Schlafst\"orungen seit 4~Monaten; CONTRA: keine Suizidalit\"at, erhaltene Arbeitsfähigkeit. \\
Ih (Verdacht): Generalisierte Angstst\"orung (F41.1), Konfidenz 55\%, PRO: GAD-7 = 12, chronische Sorgen; CONTRA: m\"oglicherweise sekund\"ar zur Depression. \\
Ii (Abdeckung): Gesamtabdeckung $\sim$78\%. Nicht abgedeckt: chronische R\"uckenschmerzen (somatisch? funktionell?), Konzentrationsprobleme (depressionsassoziiert oder eigenst\"andig?). \\
Ij (Plan): Dringend: Suizidalit\"atsmonitoring (C-SSRS); Wichtig: Neuropsychologische Testung zur Kl\"arung der Konzentrationsst\"orung; Verlauf: PHQ-9-Monitoring alle 4~Wochen.
\item[Achse~II (Biographie)]
PID-5-BF+M: Erh\"ohte Negative Affektivit\"at (T=68), grenzwertige Distanziertheit (T=60). \\
Pr\"agende Erfahrung: Fr\"uhe Scheidung der Eltern (Alter 7), intermittierendes Mobbing in der Schule (12--15~J.). \\
Grundkonflikt: Autonomie-Abh\"angigkeits-Konflikt (OPD).
\item[Achse~III (Med. Synopse)]
IIIa: Lumbales Schmerzsyndrom (M54.5). IIIc: Hypothyreose (E03.9), beitragend zur Fatigue --- Kausalanalyse (IIIk): TSH = 6.2~mU/L, partielle Erkl\"arung der Antriebsminderung. \\
IIIm: L-Thyroxin 75\,$\mu$g, Wirkung 6/10; Ibuprofen 400\,mg b.B., Wirkung 5/10; keine relevanten Interaktionen.
\item[Achse~IV (Umwelt \& Funktion)]
WHODAS~2.0 Gesamtscore: 28/100 (moderate Beeintr\"achtigung). Mini-ICF-APP: Einschr\"ankungen in Regelm\"a\ss{}igkeit (3/4), Kontaktf\"ahigkeit (2/4). GdB: nicht beantragt. \\
Stressoren: Drohender Arbeitsplatzverlust, Partnerschaftskrise. \\
Bezugspersonen: Partnerin (ambivalent st\"utzend), Hausarzt Dr.~X, kein Therapeut bisher.
\item[Achse~V (Bedingungsmodell)]
P1 (Pr\"adisposition): Fr\"uhe Verlusterfahrung, erh\"ohte Negative Affektivit\"at. \\
P2 (Ausl\"oser): Arbeitsplatzbedrohung seit 4~Monaten, zeitlich konkordant mit Symptombeginn. \\
P3 (Aufrechterhaltung): Sozialer R\"uckzug, unbehandelte Hypothyreose, Schmerzchronifizierung. \\
P4 (Protektiv): Durchschnittliche Intelligenz, stabile Wohnsituation, Problembewusstsein.
\item[Achse~VI (Belegsammlung)]
Evidenz-Matrix: PHQ-9 (21.02.2026), GAD-7 (21.02.2026), Labor (TSH, CRP --- 15.02.2026). \\
CAVE: Hypothyreose beitragende Ursache der Fatigue --- Schilddr\"usenparameter vor Antidepressiva-Einstellung kontrollieren! \\
Symptomverlauf: Antriebslosigkeit seit November 2025, progredient; R\"uckenschmerzen seit 2024, fluktuierend.
\end{description}

\noindent Diese Vignette illustriert, wie das 6-Achsen-Modell eine integrierte diagnostische Perspektive erm\"oglicht: Die Abdeckungsanalyse (78\%) identifiziert ungekl\"arte Symptome, der CAVE-Hinweis verhindert eine vorschnelle Antidepressiva-Einstellung ohne Schilddr\"usenkontrolle, und das Bedingungsmodell verbindet biographische Vulnerabilit\"at mit aktuellem Ausl\"oser.


% ===================================================================
% 15. ZUSAMMENFASSUNG
% ===================================================================
\section{Zusammenfassung}
\label{sec:zusammenfassung}

Die vorliegende Arbeit hat ein 6-Achsen-Modell zur computergest\"utzten multiaxialen psychiatrischen Diagnostik vorgestellt, das DSM-5-TR, ICD-11 und ICF in einem integrierten Expertensystem vereint. Das Modell adressiert jede historische Kritik am DSM-IV-System und f\"ugt neuartige Komponenten hinzu: die quantitative Abdeckungsanalyse mit Symptom-Diagnosen-Matrix und prozentualer Gesamtmetrik, das operationalisierte Bedingungsmodell, die PRO/CONTRA-Evidenzbewertung mit Konfidenzsch\"atzung pro Diagnose, den priorisierten 3-Stufen-Untersuchungsplan, CAVE-Warnhinweise f\"ur achsen\"ubergreifende klinische Risiken, den longitudinalen Symptomverlauf, das Kontakt- und Beobachtungsprotokoll, pr\"agende Erfahrungen und Grundkonflikte (Achse~II), das strukturierte Bezugspersonen-Netzwerk (Achse~IV) und die symmetrische Achsenarchitektur mit eigenst\"andiger Medikamentenanamnese. Die Designphilosophie des Skelett-Prinzips, des lebenden Dokuments und des Berichts als Snapshot gew\"ahrleistet maximale diagnostische Freiheit bei gleichzeitiger struktureller St\"utze. Die 6-Stufen-Gatekeeper-Logik bildet den Goldstandard nach First (2024) exakt ab. Die HiTOP-Spektren werden automatisch aus Cross-Cutting-Daten berechnet. Die exemplarische multi-professionelle Zuordnung bietet eine flexible, nicht verpflichtende Orientierung f\"ur Psychologen, Mediziner und Sozialarbeiter. Die Integration dimensionaler Perspektiven (HiTOP, RDoC), ethischer Leitlinien und einer rigoros gestuften Validierungsstrategie sichert wissenschaftliche Anschlussf\"ahigkeit. Die technische Implementierung als hierarchische Zustandsmaschine mit 11~St\"orungsmodulen und FHIR-Interoperabilit\"at bietet einen klaren Pfad vom Prototyp zum klinischen Werkzeug. Die institutionelle Validierung durch das APA Future DSM Strategic Committee und die Forderungen von Erlich und First (2025) best\"atigen die Relevanz des Ansatzes.


% ===================================================================
% LITERATURVERZEICHNIS
% ===================================================================
\newpage
\bibliographystyle{apalike}
\bibliography{diagnostic_references}

% ===================================================================
% ENGLISCHE ÜBERSETZUNG
% ===================================================================
\end{document}
